\documentclass[11pt,oneside]{article}
\usepackage[T1]{fontenc}
\usepackage[utf8]{inputenc}
\usepackage{lmodern}
\usepackage[margin=0.88in]{geometry}
\usepackage{microtype}
\usepackage{amsmath,amssymb,amsthm,mathtools}
\usepackage{booktabs,longtable,tabularx,array}
\usepackage{enumitem}
\usepackage{xcolor}
\usepackage{listings}
\usepackage{fancyhdr}
\usepackage{graphicx}
\usepackage[hyphens]{url}
\usepackage{xurl}
\usepackage[hypertexnames=false]{hyperref}
\usepackage{bookmark}

\hypersetup{
  colorlinks=true,
  linkcolor=blue!50!black,
  citecolor=blue!50!black,
  urlcolor=blue!65!black,
  pdftitle={Bring radicals and quintic equations in intuitionistic primitive recursive arithmetic},
  pdfauthor={OpenAI, prepared for Vladimir Reshetnikov},
  pdfsubject={Constructive finite proofs, exact certificates, and a Lean formalization roadmap}
}

\pagestyle{fancy}
\fancyhf{}
\lhead{Bring radicals and quintics in iPRA}
\rhead{\thepage}
\setlength{\headheight}{14pt}
\setlength{\emergencystretch}{3em}
\allowdisplaybreaks
\setlist{leftmargin=*,itemsep=3pt,topsep=5pt}

\newtheorem{theorem}{Theorem}[section]
\newtheorem{lemma}[theorem]{Lemma}
\newtheorem{proposition}[theorem]{Proposition}
\newtheorem{corollary}[theorem]{Corollary}
\newtheorem{claim}[theorem]{Claim}
\theoremstyle{definition}
\newtheorem{definition}[theorem]{Definition}
\newtheorem{construction}[theorem]{Construction}
\newtheorem{example}[theorem]{Example}
\newtheorem{audit}[theorem]{Audit item}
\theoremstyle{remark}
\newtheorem{remark}[theorem]{Remark}
\newtheorem{frontier}[theorem]{Formalization frontier}

\newcommand{\N}{\mathbb N}
\newcommand{\Z}{\mathbb Z}
\newcommand{\Q}{\mathbb Q}
\newcommand{\R}{\mathbb R}
\newcommand{\C}{\mathbb C}
\newcommand{\F}{\mathbb F}
\newcommand{\ipra}{\ensuremath{\mathsf{iPRA}}}
\newcommand{\BR}{\operatorname{BR}}
\newcommand{\Gal}{\operatorname{Gal}}
\newcommand{\Aut}{\operatorname{Aut}}
\newcommand{\Disc}{\operatorname{Disc}}
\newcommand{\Res}{\operatorname{Res}}
\newcommand{\sgn}{\operatorname{sgn}}
\newcommand{\ord}{\operatorname{ord}}
\newcommand{\coeff}[2]{[X^{#1}]\,#2}
\newcommand{\abs}[1]{\left|#1\right|}
\newcommand{\norm}[1]{\left\lVert#1\right\rVert}
\newcommand{\angles}[1]{\left\langle #1\right\rangle}
\newcommand{\code}[1]{\texttt{#1}}
\newcommand{\Exact}{\textsf{Exact}}
\newcommand{\Effective}{\textsf{Effective}}
\newcommand{\Conditional}{\textsf{Certificate}}
\newcommand{\Outline}{\textsf{Outline}}

\lstdefinestyle{code}{
  basicstyle=\ttfamily\footnotesize,
  columns=fullflexible,
  breaklines=true,
  frame=single,
  showstringspaces=false,
  keepspaces=true
}

\begin{document}

\begin{titlepage}
\centering
\vspace*{0.85in}
{\LARGE\bfseries Bring Radicals and Quintic Equations\\[0.35em]
in Intuitionistic Primitive Recursive Arithmetic\par}
\vspace{0.7in}
{\large Finite algebra, effective real roots, solvable families,\\
and the Abel--Ruffini obstruction\par}
\vspace{0.7in}
{\large Prepared for Vladimir Reshetnikov\par}
\vfill
\begin{minipage}{0.91\textwidth}
\small
This report is written against the ProveIt repository state headed by commit
\code{b73ff58ad10e9c446fa997777c5f502844249929} (16 November 2026).
It reuses the kernel-checked \ipra{} derivations of the quadratic, Cardano and
Ferrari factorization identities, the existing finite Lagrange-inversion
packet, and the repository's generic Abel--Ruffini architecture.  The exact
Python checker in this archive independently verifies the new polynomial,
finite-field and finite-group certificates.  The accompanying Lean source is a
detailed blueprint; only the imported ProveIt derivations, not every new target
in that blueprint, are claimed to be kernel-checked.
\end{minipage}
\vspace{0.55in}
{\large 16 November 2026\par}
\end{titlepage}

\pagenumbering{roman}

\begin{abstract}
We reconstruct a large part of the mathematical content of the reference pages
on the Bring radical and quintic equations in a form suitable for
intuitionistic primitive recursive arithmetic.  The governing discipline is
that a purported infinite or analytic object is replaced by a
primitive-recursive name together with an explicit modulus, while every
algebraic assertion is reduced to a finite polynomial identity, a bounded
search, or a finite certificate.

For real input we construct the Bring radical by exact rational bisection and
prove uniqueness, oddness, strict decrease, unboundedness, a sharp inverse
Lipschitz bound, a residual error bound, differentiability, and an explicit
fifth-root asymptotic estimate.  Lagrange inversion in truncated power-series
rings yields the Taylor coefficients and their hypergeometric form.  A
constructive ratio argument gives the exact convergence radius
\(4/(5\sqrt[4]{5})\); the resultant
\(3125a^4+256\) simultaneously identifies the four complex branch points.
We also verify the Cockle--Harley fourth-order differential resolvent.

For general quintics we prove the exact depression formula and a transparent
quadratic Tschirnhaus reduction to principal form.  We formulate the enormous
Bring--Jerrard step as a compact resultant certificate whose checking is
primitive recursive, and prove the scaling formula expressing
\(x^5+px+q\) through a Bring radical.  The generic Abel--Ruffini obstruction is
then rendered finitarily: the generic Galois group is \(S_5\), radical towers
have solvable normal closures, and explicit commutator certificates show that
\(S_5\) is not solvable.  Two concrete polynomials,
\(x^5-x+1\) and \(x^5+x+3\), are certified to have group \(S_5\) by
factorizations modulo small primes.  In particular \(\BR(3)\) is not a radical
number over \(\Q\).

The solvable side is developed by exact cyclotomic identities.  We prove the
de Moivre quintic formula, MathWorld's four-resolvent expansion, several
explicit radical and cyclotomic examples, and a finite classification check:
every solvable transitive subgroup of \(S_5\) lies in a conjugate of
\(\operatorname{AGL}(1,5)\).  Finally, denominator clearing is used to derive
the quintic occurring in the collinear Lagrange-point problem.  A statement
audit distinguishes completed exact arguments, effective analytic arguments,
certificate theorems, and remaining global modular-function formalization.
\end{abstract}

\tableofcontents
\clearpage
\pagenumbering{arabic}

\section{Scope, conventions, and the constructive reading}\label{sec:scope}

\subsection{The three source pages}

The requested source material consists of the Wikipedia articles
\emph{Bring radical} and \emph{Quintic equation}, and the MathWorld article
\emph{Quintic Equation}.  They mix several logically different kinds of
assertion:

\begin{enumerate}
\item finite algebraic identities, such as a Tschirnhaus substitution, a
      resultant, a cyclotomic product, or a displayed radical formula;
\item existence and approximation assertions for real or complex roots;
\item convergence, branch-point, differential-equation, theta-function, and
      hypergeometric assertions;
\item Galois-theoretic impossibility and solvability criteria;
\item historical attributions and computational advice.
\end{enumerate}

Only the first four are mathematical propositions.  Historical priority is not
an arithmetic theorem and is therefore recorded only as bibliography.  The
three pages were inspected on 16 November 2026.  Their main claims relevant to
this report are the following.

\begin{itemize}
\item The real Bring radical \(\BR(a)\) is the unique real zero of
      \(x^5+x+a\); it is odd, decreasing and unbounded, with
      \(\BR(a)\sim-a^{1/5}\).  Its local series, hypergeometric form, radius,
      and four branch points are displayed.
\item A general quintic can first be depressed, then put in principal form,
      and finally related to a Bring--Jerrard equation by a quartic
      Tschirnhaus transformation.  A trinomial \(x^5+px+q\) is reduced to one
      variable of the Bring radical by scaling.
\item General quintics are not solvable by radicals, while special quintics
      are.  The pages give de Moivre's family, resolvent constructions,
      solvable parameterizations, explicit examples, and finite-group
      criteria.
\item Beyond radicals, the pages describe hypergeometric, theta, elliptic,
      icosahedral and iterative methods, and the quintic arising from the
      collinear Lagrange points.
\end{itemize}

The goal is not to copy those pages, but to identify which claims admit a clean
finitary proof and to expose the exact data needed by \ipra.

\subsection{Four meanings of ``proved in \(\ipra\)''}

First-order \ipra{} quantifies only over natural numbers.  It has no primitive
sort of arbitrary real numbers, complex functions, fields, or infinite
series.  We therefore use four precise readings.

\begin{definition}[Exact finite statement]
An \emph{exact} statement is an identity or decidable relation on natural-number
codes for integers, rationals, finite vectors, polynomials, permutations, or
finite tables.  Denominators are cleared and all loops have explicit bounds.
Such a statement is directly expressible in \ipra.
\end{definition}

\begin{definition}[Effective analytic statement]
A real or complex quantity is represented by a fixed primitive-recursive
sequence of rational or Gaussian-rational approximants, together with a fixed
primitive-recursive modulus.  An analytic theorem is a theorem schema for such
named data, with all radii, bounds and error moduli explicit.  No quantifier
ranges over an unspecified completed function space.
\end{definition}

\begin{definition}[Certificate statement]
A certificate statement says that a primitive-recursive verifier accepts a
finite object only when a specified mathematical relation holds.  Examples are
a resultant identity, a factorization modulo a prime, a multiplication table
for a finite Galois group, or a Tschirnhaus transformation.  The theorem is
uniform in the certificate code; a concrete certificate is then checked by
finite computation.
\end{definition}

\begin{definition}[Formalization outline]
An outline gives all object-theory definitions, the exact lemmas to be proved,
and a dependency path to a kernel-checked Lean theorem, but does not claim that
the final Lean term has already been compiled.  We use this label for the
large global theta-function and universal Bring--Jerrard developments.
\end{definition}

The audit in Section~\ref{sec:audit} labels results by
\(\Exact\), \(\Effective\), \(\Conditional\), or \(\Outline\).  These labels
are deliberately stronger than a vague assertion that an ordinary classical
proof is ``finitary in spirit.''

\subsection{Sign conventions}

Two closely related normal forms occur in the sources.  Throughout this report
we define
\begin{equation}\label{eq:BRdef}
  \BR(a)^5+\BR(a)+a=0
\end{equation}
for real \(a\).  Thus if
\[
  h(x)=x^5+x,
\]
then \(\BR(a)=h^{-1}(-a)=-h^{-1}(a)\).  The alternative Bring equation
\(t^5-t+\rho=0\) is treated separately; it has different critical values and
is the convention used in the Cockle--Harley differential resolvent.

For a primitive fifth root of unity we write \(\omega\), so
\begin{equation}\label{eq:cyclo5}
 \omega^5=1,
 \qquad 1+\omega+\omega^2+\omega^3+\omega^4=0.
\end{equation}
No analytic choice of \(e^{2\pi i/5}\) is needed for the algebraic identities:
\(\omega\) is simply an element satisfying these two equations.

\subsection{What is new in this packet}

Several ingredients already existed in ProveIt.  The repository's
\code{IPRADerivations} library contains kernel-checked object-theory proofs of
the quadratic, Cardano and Ferrari factorization identities.  Its Lagrange
inversion packet gives the finite coefficient algebra needed for inverses of
\(x-x^p\).  A separate polynomial-formulas packet develops a generic
Abel--Ruffini architecture.  We use these rather than silently reproving them
in a stronger metatheory.

The new contribution is the synthesis specific to Bring radicals and the
quintic pages:

\begin{itemize}
\item a short sum-of-squares proof of strong monotonicity and the resulting
      primitive-recursive real Bring radical;
\item exact coefficient and branch-point proofs tied to the finite Lagrange
      framework;
\item a transparent principal-form reduction and a compact certificate
      interface for the large Bring--Jerrard step;
\item concrete finite-field certificates proving that both
      \(x^5-x+1\) and \(x^5+x+3\) have Galois group \(S_5\);
\item exact de Moivre and four-resolvent identities, plus a complete finite
      enumeration of the subgroup classification relevant to solvable
      irreducible quintics;
\item a detailed Lean module plan connected to the existing ProveIt theorem
      names.
\end{itemize}

The executable verifier performs seventeen independent families of exact
checks.  It is evidence for the displayed certificates, not a substitute for
the human proofs or the intended object-theory derivations.

\section{The finite infrastructure inside \(\ipra\)}\label{sec:framework}

\subsection{The object theory actually used}

The ProveIt presentation of \ipra{} is first-order intuitionistic arithmetic
with equality, a symbol and defining equations for every fixed primitive
recursive function, the successor axioms, and induction for quantifier-free
formulas.  The crucial word is \emph{fixed}: we may add the coefficient
function of one series, the verifier of one certificate format, or the step
function of one iteration.  We do not postulate a primitive-recursive universal
evaluator for all primitive-recursive programs.

Every invariant below is represented by an equality of natural-number terms.
For example, a bounded assertion
\[
  \forall i<N\; P(i)
\]
with decidable \(P\) is encoded by the Boolean fold
\[
  \operatorname{All}_P(0)=1,
  \qquad
  \operatorname{All}_P(i+1)=
  \operatorname{All}_P(i)\cdot\chi_P(i).
\]
The assertion is the quantifier-free equation
\(\operatorname{All}_P(N)=1\).  Consequently the ordinary proof ``maintain
this finite invariant for every stage'' is literally quantifier-free
induction.

\subsection{Integers, rationals, and finite algebra}

An integer is coded by a pair \((u,v)\) representing \(u-v\).  Equality is
cross-addition,
\[
 (u,v)=(u',v')\quad\Longleftrightarrow\quad u+v'=u'+v.
\]
A rational is represented by a signed numerator and a positive denominator;
all rational identities are converted to integer identities by multiplying by
the product of the guarded nonzero denominators.  The ProveIt
\code{IPRADerivations.Core.Integers} and \code{Core.Tactic} layers already
implement this difference-ring method and a ring normalizer for object-theory
polynomial identities.

A polynomial of degree at most \(d\) is a finite coefficient vector.  Addition,
multiplication, formal differentiation, evaluation, Euclidean division over a
field, Sylvester determinants and resultants are bounded loops.  A sparse
representation is convenient for the very large universal Tschirnhaus
expressions, but no theorem depends on the choice of representation.

\begin{lemma}[Finite polynomial normalization]\label{lem:poly-normalize}
For each fixed number of variables and each coded degree bound there is a
primitive-recursive normalizer \(\operatorname{nf}\) for integer polynomials
such that \ipra{} proves
\[
 \operatorname{nf}(P)=\operatorname{nf}(Q)
 \quad\Longrightarrow\quad
 P(\mathbf x)=Q(\mathbf x).
\]
The same is true for rational polynomials after denominator clearing.
\end{lemma}

\begin{proof}
Encode a monomial by its exponent vector and a polynomial by a finite list of
coefficient--monomial pairs.  Insertion sort by the lexicographic monomial code,
followed by a bounded fold combining equal monomials, is primitive recursive.
Correctness is proved by induction on the input length.  The invariant is an
equality between the evaluation of the processed prefix and the evaluation of
the normalized accumulator, hence quantifier-free.  Rational coefficients are
handled by a common positive denominator.
\end{proof}

This elementary lemma is the object-theory meaning of every use of
\code{ring}, \code{field\_simp}, or a symbolic resultant in the metatheoretic
Lean and Python artifacts.

\subsection{Regular names for real and complex numbers}

A \emph{regular dyadic name} for a real number is a primitive-recursive
function \(q:\N\to\Q\) satisfying
\begin{equation}\label{eq:regular-name}
  \abs{q(n)-q(m)}\le 2^{-n}+2^{-m}.
\end{equation}
Two names are equal if
\(\abs{q(n)-r(n)}\le 2^{1-n}\) for all \(n\).  A complex name is a pair of real
names.  Arithmetic on names uses explicit error propagation.  On a supplied
rational bound \(\abs{x}\le M\), the polynomial map \(x\mapsto x^5+x\) has an
explicit rational modulus because
\[
 \abs{x^5-y^5}
 \le 5(M+1)^4\abs{x-y}.
\]
after increasing the bound by one to absorb approximation error.

Statements about order use located intervals rather than an illicit global law
of trichotomy for arbitrary reals.  In the Bring construction all branching is
performed on exact rational midpoint evaluations.  The extension from rational
parameters to a named real parameter is then made by the Lipschitz estimate of
Theorem~\ref{thm:bring-lipschitz}; no undecidable sign test on a real number is
required.

\begin{definition}[Effective convergence]
A sequence \(z_n\) converges to \(z\) with modulus \(M\) when
\[
 m,n\ge M(k)\quad\Longrightarrow\quad
 \abs{z_m-z_n}\le 2^{-k}.
\]
every term and \(M\) being primitive recursive.  A power series is asserted to
converge only together with a primitive-recursive tail bound on each named
closed subdisk of its open disk of convergence.
\end{definition}

This is the reading used for the Bring series, hypergeometric functions and
local analytic branches.

\subsection{Truncated power series}

For a commutative coded ring \(R\), define
\[
  T_N(R)=R[X]/(X^{N+1})
\]
as vectors \((a_0,\ldots,a_N)\).  Cauchy multiplication is
\[
  (ab)_n=\sum_{i=0}^{n}a_i b_{n-i}.
\]
Composition \(A(B)\) is a bounded sum whenever \(B_0=0\).  If \(U_0=1\), the
inverse of \(U\) in \(T_N(R)\) is produced by the finite geometric series in
\(1-U\), or recursively coefficient by coefficient.  All these operations are
primitive recursive in \(N\).

The finite Lagrange theorem used later can be stated without completed series.

\begin{theorem}[Finite Lagrange inversion, reusable form]\label{thm:finite-lagrange}
Let \(\Phi\in T_N(\Q)\).  Starting from \(G_0=0\), define
\[
 G_{r+1}=X\Phi(G_r)\pmod{X^{N+1}}.
\]
Then \(G_{N+1}\) is the unique vector with zero constant term satisfying
\(G=X\Phi(G)\).  For \(1\le n\le N\),
\begin{equation}\label{eq:finite-lagrange}
 n[X^n]G=[W^{n-1}]\Phi(W)^n.
\end{equation}
Every operation and the proof of the theorem are primitive recursive in \(N\),
and the theorem is derivable in \ipra{} by quantifier-free induction.
\end{theorem}

\begin{proof}
If two vectors agree below degree \(m\), their images under
\(A\mapsto X\Phi(A)\) agree below degree \(m+1\).  Thus the successive iterates
stabilize one coefficient per step, which gives existence and uniqueness after
\(N+1\) steps.  Formula~\eqref{eq:finite-lagrange} follows from the finite
Laurent-residue substitution identity
\[
 \Res Q(F(W))F'(W)=\Res Q(W).
\]
after truncating every Laurent vector to the finite exponent interval that can
contribute to the residue.  The proof on Laurent monomials is a telescoping
coefficient identity; linearity gives the general case.  This is exactly the
architecture of the existing ProveIt Lagrange-inversion packet.
\end{proof}

Writing the equation with the factor \(n\) on the left avoids division in the
object theory.  Division by \(n>0\) is used only in the rational metanotation.

\subsection{Finite fields, permutations, and groups}

For the Galois arguments we use two complementary encodings.

\begin{enumerate}
\item For a concrete polynomial modulo a concrete prime, elements of
      \(\F_p[X]/(f)\) are coefficient vectors of length \(\deg f\).  Addition,
      multiplication, gcd, Frobenius powers and Bézout certificates are
      bounded algorithms.
\item An abstract finite group is a natural number \(n\) together with an
      \(n\times n\) multiplication table and tables for identity and inverse.
      The group axioms, subgroup closure, commutators, derived series, orbits
      and permutation actions are bounded Boolean folds.
\end{enumerate}

A permutation of five letters is a five-entry vector.  There are only
\(5!=120\) candidates.  The executable certificate enumerates all 156 actual
subgroups of \(S_5\), identifies the 20 transitive ones, and checks that the 18
solvable transitive subgroups have orders \(5,10\), or \(20\) and are contained
in a conjugate of \(\operatorname{AGL}(1,5)\).  This is a finite computation,
so its trace can be reproduced by a primitive-recursive verifier inside
\ipra.

\begin{lemma}[Finite certificate reflection]\label{lem:certificate-reflection}
Let \(V(c)\in\{0,1\}\) be a fixed primitive-recursive verifier, and suppose
\ipra{} proves the uniform soundness implication
\[
 V(c)=1\longrightarrow R(c).
\]
after all bounded assertions in \(R\) have been Boolean-folded.  Then every
explicit accepted certificate \(\bar c\) yields an \ipra{} proof of
\(R(\bar c)\).
\end{lemma}

\begin{proof}
The closed computation \(V(\bar c)=1\) reduces by the defining equations of the
primitive-recursive symbols.  Modus ponens with the uniform soundness theorem
gives \(R(\bar c)\).  No reflection principle for arbitrary proofs is being
assumed; this is only evaluation of one fixed decidable verifier.
\end{proof}

\subsection{Reuse of the checked ProveIt library}

The following dependencies are already available in the repository.

\begin{itemize}
\item \code{IPRADerivations.Core.Semiring}, \code{Core.Integers},
      \code{Core.Order}, \code{Core.Powers}, \code{Core.Search} and
      \code{Core.Pairing} provide the basic arithmetic, bounded search and
      coding infrastructure.
\item \code{IPRADerivations.Algebra.PolynomialFormulas.quadratic},
      \code{cardano}, and \code{ferrari} are closed derivations of the exact
      factorization identities needed to solve the auxiliary equations of
      degrees two, three and four.
\item The transfer theorem \code{pra\_iff\_native} transports these ordinary
      \ipra{} derivations to the repository's very small native scalar core.
\end{itemize}

The present Lean blueprint is designed to add new object-theory modules above
this checked base rather than to replace it with Mathlib field reasoning.

\section{The real Bring radical as a primitive-recursive real}\label{sec:real-bring}

Set
\[
 P_a(x)=x^5+x+a,
 \qquad h(x)=x^5+x.
\]
The usual proof that \(h\) is strictly increasing appeals to ordinary real
calculus.  For \ipra{} it is better to use a polynomial certificate that is
valid over every ordered field.

\subsection{A sum-of-squares monotonicity certificate}

Define
\[
 S_4(x,y)=x^4+x^3y+x^2y^2+xy^3+y^4.
\]

\begin{lemma}[Positivity of the fifth-power divided difference]
\label{lem:s4-sos}
The following identity holds in \(\Z[x,y]\):
\begin{equation}\label{eq:s4-sos}
 12S_4(x,y)
 =3(2x^2+xy)^2+(3xy+2y^2)^2+8y^4.
\end{equation}
Consequently \(S_4(x,y)\ge0\) in every ordered field.
\end{lemma}

\begin{proof}
Expansion of the right side gives
\[
 12x^4+12x^3y+12x^2y^2+12xy^3+12y^4.
\]
This is an exact polynomial normalization, so after coding signed coefficients
it is an \ipra{} ring identity.  Nonnegativity follows from nonnegativity of
squares and multiplication by the positive integer \(12\).
\end{proof}

\begin{theorem}[Strong monotonicity]\label{thm:strong-monotonicity}
For \(y\ge x\),
\begin{equation}\label{eq:strong-monotonicity}
 P_a(y)-P_a(x)
 =(y-x)(1+S_4(x,y))\ge y-x.
\end{equation}
In particular \(P_a\) and \(h\) are strictly increasing.
\end{theorem}

\begin{proof}
The factorization
\(y^5-x^5=(y-x)S_4(x,y)\) is exact.  Apply
Lemma~\ref{lem:s4-sos}.  If \(y>x\), then the right side is at least the
positive number \(y-x\).
\end{proof}

The extra lower bound in~\eqref{eq:strong-monotonicity}, rather than mere strict
increase, is the key constructive estimate used throughout this section.

\subsection{Exact rational bisection}

Let \(a\in\Q\).  Put
\[
 M=\lceil |a|\rceil+1.
\]
Then
\[
 P_a(-M)<0<P_a(M).
\]
Starting with \([L_0,U_0]=[-M,M]\), define
\(m_n=(L_n+U_n)/2\) and
\begin{equation}\label{eq:bisection}
 (L_{n+1},U_{n+1})=
 \begin{cases}
   (m_n,U_n),&P_a(m_n)\le0,\\
   (L_n,m_n),&P_a(m_n)>0.
 \end{cases}
\end{equation}
Every comparison in~\eqref{eq:bisection} is a decidable comparison of rational
codes.

\begin{lemma}[Bisection invariant]\label{lem:bisection-invariant}
For every \(n\),
\[
 L_n\le U_n,
 \qquad P_a(L_n)\le0\le P_a(U_n),
 \qquad U_n-L_n=2M\,2^{-n}.
\]
The functions \(L_n,U_n\) are primitive recursive in the code of \(a\) and
\(n\).
\end{lemma}

\begin{proof}
One step of~\eqref{eq:bisection} preserves the two endpoint signs by the branch
condition and halves the width.  All statements are rational equalities or
orders and the induction invariant can be Boolean-folded, so quantifier-free
induction suffices.
\end{proof}

Let
\[
 N(a,k)=k+2+\left\lceil\log_2(2M)\right\rceil,
 \qquad
 q_a(k)=\frac{L_{N(a,k)}+U_{N(a,k)}}2.
\]
Then the containing interval has width at most \(2^{-k-1}\).  The sequence
\(q_a\) is therefore a regular dyadic name.

\begin{theorem}[Constructive existence and uniqueness]\label{thm:bring-exists}
For every rational \(a\), the name \(q_a\) converges with its displayed
primitive-recursive modulus to a real \(x\) satisfying
\[
 x^5+x+a=0.
\]
This zero is unique.
\end{theorem}

\begin{proof}
Nested intervals with a geometric width give a regular Cauchy name.  On
\([-M,M]\), the exact factorization of fifth powers gives the rational
Lipschitz estimate
\[
 |P_a(u)-P_a(v)|
 \le \bigl(1+5M^4\bigr)|u-v|.
\]
Since \(P_a(L_n)\le0\le P_a(U_n)\), monotonicity places
\(P_a(q_a(k))\) between the endpoint values; the difference of those endpoint
values tends effectively to zero by the displayed Lipschitz bound.  Hence the
limit satisfies \(P_a(x)=0\).  If \(x<y\) were two zeros,
Theorem~\ref{thm:strong-monotonicity} would give
\(0=P_a(y)-P_a(x)\ge y-x>0\), a contradiction.
\end{proof}

\begin{definition}[Real Bring radical]
For rational \(a\), \(\BR(a)\) is the regular name constructed above.  For a
regular real name \(a_n\), define \(\BR(a)\) by applying rational bisection to
a sufficiently accurate \(a_n\) and using the stability estimate in the next
theorem.  This extension is independent of the chosen name.
\end{definition}

\subsection{Stability, oddness, and order}

\begin{theorem}[Sharp inverse Lipschitz estimate]\label{thm:bring-lipschitz}
For all real \(a,b\),
\begin{equation}\label{eq:bring-lipschitz}
 \abs{\BR(a)-\BR(b)}\le |a-b|.
\end{equation}
More generally, for every real \(q\),
\begin{equation}\label{eq:residual-bound}
 \abs{q-\BR(a)}\le |q^5+q+a|.
\end{equation}
\end{theorem}

\begin{proof}
Let \(x=\BR(a)\), \(y=\BR(b)\), and suppose \(x\le y\).  Strong monotonicity
gives
\[
 y-x\le h(y)-h(x)=(-b)-(-a)=a-b\le|a-b|.
\]
The other order is symmetric.  Taking \(y=q\) and using
\(P_a(x)=0\) gives~\eqref{eq:residual-bound}.
\end{proof}

Estimate~\eqref{eq:bring-lipschitz} proves at once that applying rational
bisection to successively better rational approximants of a named real gives a
regular name.  This is the constructive extension promised above.

\begin{corollary}[Oddness and strict decrease]\label{cor:bring-odd-decreasing}
The Bring radical is odd and strictly decreasing:
\[
 \BR(-a)=-\BR(a),
 \qquad a<b\Longrightarrow \BR(b)<\BR(a).
\]
\end{corollary}

\begin{proof}
If \(x^5+x+a=0\), then \((-x)^5+(-x)-a=0\); uniqueness gives oddness.  If
\(a<b\), then
\[
 P_b(\BR(a))=b-a>0.
\]
Since \(P_b\) is strictly increasing and vanishes at \(\BR(b)\), its positive
value at \(\BR(a)\) forces \(\BR(b)<\BR(a)\).
\end{proof}

\begin{corollary}[Unboundedness with a witness]\label{cor:bring-unbounded}
For every \(m\ge0\),
\[
 a>m^5+m\Longrightarrow \BR(a)<-m,
 \qquad
 a<-(m^5+m)\Longrightarrow \BR(a)>m.
\]
Thus \(\BR\) is unbounded in both directions.
\end{corollary}

\begin{proof}
For the first implication,
\(P_a(-m)=a-m^5-m>0\).  The zero of the increasing polynomial lies to the
left of \(-m\).  The second implication follows by oddness.
\end{proof}

\subsection{Differentiability without an inverse-function black box}

Let \(x=\BR(a)\), \(y=\BR(a+t)\).  Subtracting the two defining equations gives
the exact identity
\begin{equation}\label{eq:bring-difference-quotient}
 (y-x)\bigl(1+S_4(x,y)\bigr)=-t.
\end{equation}
The denominator is at least one by Lemma~\ref{lem:s4-sos}.

\begin{theorem}[Derivative of the real Bring radical]\label{thm:bring-derivative}
The real Bring radical is effectively differentiable and
\begin{equation}\label{eq:bring-derivative}
 \BR'(a)=-\frac{1}{1+5\BR(a)^4}.
\end{equation}
On every rationally bounded interval, the convergence of the difference
quotient has a primitive-recursive modulus.
\end{theorem}

\begin{proof}
By~\eqref{eq:bring-lipschitz}, \(|y-x|\le|t|\).  Hence
\(S_4(x,y)\to S_4(x,x)=5x^4\) with an explicit polynomial modulus on any
supplied bound for \(|x|\).  Divide~\eqref{eq:bring-difference-quotient} by
\(t\ne0\) and pass to the effective limit.  Since the denominator is uniformly
at least one, reciprocal error propagation is primitive recursive.
\end{proof}

Repeated differentiation produces rational functions of \(\BR(a)\) whose
denominators are powers of \(1+5\BR(a)^4\).  For example,
\[
 \BR''(a)=-\frac{20\BR(a)^3}{(1+5\BR(a)^4)^3}.
\]
Thus each fixed derivative is a named continuous function.  A quantitative
complex inverse theorem applied at a real point gives a local analytic branch;
all constants are computable because \(5x^4+1\ge1\) on the real axis.  This
establishes, in the effective sense, the source page's assertion that the
complex branch can be chosen analytic in a neighborhood of the real line.

\subsection{Fifth-root asymptotics}

For \(a>0\), put \(y=-\BR(a)>0\).  Then
\begin{equation}\label{eq:y-equation}
 y^5+y=a.
\end{equation}
Let \(r=a^{1/5}\), represented by its standard primitive-recursive bisection
name.  Equation~\eqref{eq:y-equation} gives \(y\le r\).  Factoring the
difference of fifth powers yields
\begin{equation}\label{eq:asymptotic-exact}
 (r-y)(r^4+r^3y+r^2y^2+ry^3+y^4)=y.
\end{equation}

\begin{theorem}[Effective asymptotic bound]\label{thm:bring-asymptotic}
For \(a=r^5\) with \(r\ge1\),
\begin{equation}\label{eq:bring-asymptotic-bound}
 0\le r+\BR(a)\le r^{-3},
 \qquad
 0\le 1-\frac{-\BR(a)}{r}\le r^{-4}.
\end{equation}
Consequently \(\BR(a)\sim-a^{1/5}\) as \(a\to+\infty\), with the explicit
modulus supplied by~\eqref{eq:bring-asymptotic-bound}.  The negative direction
follows by oddness.
\end{theorem}

\begin{proof}
The parenthesized factor in~\eqref{eq:asymptotic-exact} is at least \(r^4\),
and \(y\le r\).  Therefore
\(0\le r-y\le y/r^4\le r^{-3}\).  Divide by \(r>0\) for the relative bound.
All inequalities become integer polynomial inequalities after multiplying by
positive powers of \(r\), so the proof is compatible with the regular-real
coding.
\end{proof}

\section{Series, hypergeometric form, branch points, and the differential resolvent}\label{sec:series}

\subsection{The Taylor coefficients by finite Lagrange inversion}

Put
\[
 y(a)=-\BR(a).
\]
Then \(y(0)=0\) and
\begin{equation}\label{eq:y-fixed}
 y+y^5=a,
 \qquad
 y=a(1+y^4)^{-1}.
\end{equation}
In the truncated ring \(\Q[a]/(a^{N+1})\), the inverse
\((1+y^4)^{-1}\) is a finite geometric sum.  Theorem~\ref{thm:finite-lagrange}
therefore applies with \(\Phi(w)=(1+w^4)^{-1}\).

\begin{theorem}[Bring coefficient formula]\label{thm:bring-coeff}
For every \(N\), the unique truncated solution of~\eqref{eq:y-fixed} is
\begin{equation}\label{eq:bring-series-y}
 y(a)=\sum_{4k+1\le N}
 (-1)^k\frac{1}{4k+1}\binom{5k}{k}a^{4k+1}
 \pmod {a^{N+1}}.
\end{equation}
Consequently
\begin{equation}\label{eq:bring-series}
 \BR(a)=\sum_{k\ge0}
 (-1)^{k+1}\frac{1}{4k+1}\binom{5k}{k}a^{4k+1}
 =-a+a^5-5a^9+35a^{13}-285a^{17}+\cdots .
\end{equation}
The assertion through every requested degree is a theorem of \ipra.
\end{theorem}

\begin{proof}
Finite Lagrange inversion gives
\[
 [a^n]y=\frac1n[w^{n-1}](1+w^4)^{-n}.
\]
The finite negative-binomial identity is
\[
 (1+w^4)^{-n}
 =\sum_{j=0}^{\lfloor N/4\rfloor}
   (-1)^j\binom{n+j-1}{j}w^{4j}
   \pmod {w^{N+1}}.
\]
Thus the coefficient vanishes unless \(n-1=4k\).  For
\(n=4k+1\),
\[
 \binom{n+k-1}{k}=\binom{5k}{k},
\]
which gives~\eqref{eq:bring-series-y}.  In the object theory one keeps the
integer equation
\((4k+1)[a^{4k+1}]y=(-1)^k\binom{5k}{k}\), so no unguarded division is used.
\end{proof}

The archive's exact checker forms the right side through degree 65 and verifies
that \(y^5+y-a\) has no coefficient through that degree.  This finite check is
not the proof of the uniform theorem, but it is a useful independent regression
test.

\subsection{The hypergeometric identity is a coefficient recurrence}

Let
\[
 c_k=\frac{1}{4k+1}\binom{5k}{k}
     =\frac{(5k)!}{k!(4k+1)!}.
\]
A direct cancellation of factorials gives
\begin{equation}\label{eq:bring-ratio}
 \frac{c_{k+1}}{c_k}
 =\frac{3125}{256}
 \frac{(k+\frac15)(k+\frac25)(k+\frac35)(k+\frac45)}
 {(k+\frac12)(k+\frac34)(k+\frac54)(k+1)}.
\end{equation}

\begin{theorem}[Hypergeometric form]\label{thm:bring-hypergeom}
On its disk of convergence,
\begin{equation}\label{eq:bring-hypergeom}
 \BR(a)=-a\,{}_4F_3\!\left(
 \begin{matrix}\frac15,\frac25,\frac35,\frac45\\
               \frac12,\frac34,\frac54
 \end{matrix};-\frac{3125}{256}a^4\right).
\end{equation}
For every truncation order, this is an exact identity of rational coefficient
vectors and is therefore an \ipra{} theorem.
\end{theorem}

\begin{proof}
The coefficient of \(a^{4k+1}\) on the right is
\[
 -\frac{(\frac15)_k(\frac25)_k(\frac35)_k(\frac45)_k}
 {(\frac12)_k(\frac34)_k(\frac54)_k k!}
 \left(-\frac{3125}{256}\right)^k.
\]
At \(k=0\) it equals \(-1\).  Its ratio to the preceding coefficient is the
negative of~\eqref{eq:bring-ratio}, exactly matching
\((-1)^{k+1}c_k\).  Equality follows by induction on \(k\).  After clearing the
fixed denominators \(2,3,4,5\), the recurrence is an integer polynomial
identity, so quantifier-free induction proves the finite vector equality.
\end{proof}

This proof is preferable in the present setting to invoking a general theory
of hypergeometric functions: the named special function is simply the regular
sequence of its rational partial sums plus the tail modulus established next.

\subsection{Exact convergence radius and boundary convergence}

Taking \(k\to\infty\) in~\eqref{eq:bring-ratio} gives
\[
 \lim_{k\to\infty}\frac{c_{k+1}}{c_k}=\frac{3125}{256}.
\]
The limit is constructive.  Subtracting \(3125/256\) from the rational function
in~\eqref{eq:bring-ratio} gives a quotient whose numerator has degree at most
three and whose positive denominator has degree four.  A fixed integer
\(C\) therefore satisfies
\begin{equation}\label{eq:ratio-modulus}
 \left|\frac{c_{k+1}}{c_k}-\frac{3125}{256}\right|
 \le\frac{C}{k+1}.
\end{equation}
The inequality is checked after multiplication by the positive denominator,
and is thus available in \ipra{} with an explicit cutoff.

\begin{theorem}[Radius of the Bring series]\label{thm:bring-radius}
The radius of convergence of~\eqref{eq:bring-series} is
\begin{equation}\label{eq:bring-radius}
 R=\left(\frac{256}{3125}\right)^{1/4}
   =\frac{4}{5\sqrt[4]{5}}.
\end{equation}
On every named closed disk \(|a|\le r<R\), the series has a
primitive-recursive geometric tail modulus.  If \(|a|>R\), its terms do not
tend to zero.  The series also converges absolutely on \(|a|=R\), with a
primitive-recursive tail modulus of order \(N^{-1/2}\).
\end{theorem}

\begin{proof}
Regard the series after its leading factor \(a\) as a series in \(t=a^4\).
Given rational \(r<R\), choose a rational \(\theta<1\) strictly above
\((3125/256)r^4\).  Estimate~\eqref{eq:ratio-modulus} supplies an explicit
\(K\) after which consecutive absolute terms have ratio at most \(\theta\),
so a geometric tail bound follows.  For \(|a|>R\), choose a rational ratio
strictly greater than one below the eventual term ratio; the terms cannot tend
to zero.

At the boundary, the finite Wallis--Stirling inequalities used in the ProveIt
Lagrange-inversion packet give a rational constant \(C_0\) such that
\[
 c_k\left(\frac{256}{3125}\right)^k
 \le \frac{C_0}{(k+1)^{3/2}}.
\]
The proof uses only factorial products, rational upper and lower bounds for
\(\log(1+u)\), and dyadic grouping.  The tail of
\(\sum(k+1)^{-3/2}\) is bounded by a constant times \(N^{-1/2}\), again by
dyadic grouping.  Every cutoff and rational constant is primitive recursive.
\end{proof}

\subsection{The four branch points and the discriminant}

A multiple root of \(x^5+x+a\) satisfies
\[
 x^5+x+a=0,
 \qquad 5x^4+1=0.
\]
The second equation gives \(x^4=-1/5\), and hence
\(x^5=-x/5\) and \(a=-4x/5\).  Therefore
\begin{equation}\label{eq:branch-equation}
 3125a^4+256=0.
\end{equation}
Conversely every solution of~\eqref{eq:branch-equation}, with
\(x=-5a/4\), satisfies both equations.

\begin{theorem}[Trinomial discriminant]\label{thm:trinomial-discriminant}
For indeterminates \(p,q\),
\begin{equation}\label{eq:trinomial-discriminant}
 \Disc_X(X^5+pX+q)=256p^5+3125q^4.
\end{equation}
In particular, the four critical values of the Bring projection are the four
roots of~\eqref{eq:branch-equation}; all have modulus \(R\) from
\eqref{eq:bring-radius}.
\end{theorem}

\begin{proof}
The discriminant is the resultant of
\(X^5+pX+q\) and \(5X^4+p\), since
\((-1)^{5\cdot4/2}=1\).  The determinant of the fixed \(9\times9\) Sylvester
matrix normalizes to the right side of~\eqref{eq:trinomial-discriminant}.
Equivalently, eliminate \(X\): a common root has
\(X^4=-p/5\) and
\(q=-4pX/5\); raising the latter equation to the fourth power and substituting
the former yields the displayed relation.  The converse is checked after
adjoining a fourth root.  Both versions are finite polynomial arguments.
\end{proof}

At a critical root \(x_0\), one has \(x_0\ne0\) and
\(P_{xx}(x_0)=20x_0^3\ne0\).  Thus the local inverse has two Puiseux branches
with leading term proportional to \((a-a_0)^{1/2}\).  Analytic continuation
around a small rational polygonal loop enclosing \(a_0\) interchanges these
two branches.  A single-valued analytic branch on the whole complex plane is
therefore impossible, and any domain for a branch agreeing with the real Bring
radical must introduce cuts.  The local square-root construction is finite in
every truncation and fits the ramified Lagrange theorem of the reused packet;
the loop interchange can be certified by a finite winding-number computation.

The critical values lie exactly on the circle \(|a|=R\).  Hence the branch
point calculation independently proves that the Taylor disk at zero cannot be
larger than the radius found by the coefficient ratio.

\subsection{The Cockle--Harley differential resolvent}

The differential-resolvent convention uses a branch \(t=t(\rho)\) satisfying
\begin{equation}\label{eq:minus-bring}
 t^5-t+\rho=0.
\end{equation}
Write \(D=5t^4-1\).  Implicit differentiation gives
\(t'=-D^{-1}\).  Repeatedly applying the total derivative
\[
 \mathcal D=\frac{\partial}{\partial\rho}
            -\frac1D\frac{\partial}{\partial t}
\]
expresses \(t'',t^{(3)},t^{(4)}\) as rational functions of \(t\), with the
fourth derivative having denominator \(D^7\).

\begin{theorem}[Fourth-order differential resolvent]\label{thm:differential-resolvent}
Every simple branch of~\eqref{eq:minus-bring} satisfies
\begin{equation}\label{eq:cockle-harley}
 \frac{256-3125\rho^4}{1155}\,t^{(4)}
 -\frac{6250\rho^3}{231}\,t^{(3)}
 -\frac{4875\rho^2}{77}\,t''
 -\frac{2125\rho}{77}\,t'
 +t=0.
\end{equation}
For every fixed Taylor order this is an \ipra{} theorem about rational
coefficient vectors; on a named disk avoiding the critical values it is an
effective analytic theorem.
\end{theorem}

\begin{proof}
Substitute the four rational derivative expressions into the left side,
multiply by \(1155D^7\), and replace \(\rho\) by \(t-t^5\) using
\eqref{eq:minus-bring}.  Polynomial normalization gives zero.  This is the
exact certificate implemented by \code{check\_differential\_resolvent} in the
archive.  In \ipra{}, the same calculation is an integer ring identity under
the guarded hypothesis \(D\ne0\).  At a simple branch the guard follows from
the nonvanishing derivative.
\end{proof}

At \(\rho=0\), the five roots are \(0,1,-1,i,-i\), all simple.  The branch at
zero is the hypergeometric branch of Theorem~\ref{thm:bring-hypergeom} after the
appropriate sign change.  For a fourth root of unity \(\zeta\), write
\(t=\zeta+u\); then
\[
 u=-\rho\bigl(4+10\zeta^3u+10\zeta^2u^2+5\zeta u^3+u^4\bigr)^{-1}.
\]
Finite fixed-point iteration constructs the other four Taylor branches.  They
all satisfy~\eqref{eq:cockle-harley}; their sum is zero by the missing quartic
coefficient, explaining why a fourth-order equation suffices for five
algebraic branches.  This gives a constructive derivation of all five local
roots even when one does not adopt the particular hypergeometric basis printed
by MathWorld.

\section{Quintic normal forms and exact Tschirnhaus certificates}\label{sec:normal-forms}

\subsection{Depressing the general quintic}

Consider
\begin{equation}\label{eq:general-quintic}
 A x^5+B x^4+C x^3+D x^2+E x+F=0,
 \qquad A\ne0.
\end{equation}
Set \(x=y-B/(5A)\) and divide by \(A\).  Direct expansion gives a depressed
quintic
\begin{equation}\label{eq:depressed-quintic}
 y^5+p y^3+q y^2+r y+s=0,
\end{equation}
where
\begin{align}
 p&=\frac{5AC-2B^2}{5A^2},\label{eq:depress-p}\\
 q&=\frac{25A^2D-15ABC+4B^3}{25A^3},\label{eq:depress-q}\\
 r&=\frac{125A^3E-50A^2BD+15AB^2C-3B^4}{125A^4},\label{eq:depress-r}\\
 s&=\frac{3125A^4F-625A^3BE+125A^2B^2D-25AB^3C+4B^5}
          {3125A^5}.\label{eq:depress-s}
\end{align}

\begin{theorem}[Exact depression identity]\label{thm:depress-quintic}
Equations~\eqref{eq:depress-p}--\eqref{eq:depress-s} satisfy the identity
\begin{equation}\label{eq:depress-identity}
 \frac1A P\!\left(y-\frac{B}{5A}\right)
 =y^5+py^3+qy^2+ry+s.
\end{equation}
After multiplication by \(3125A^5\), this is an integer polynomial identity
and hence a theorem of \ipra.
\end{theorem}

\begin{proof}
Expand the left side, collect powers of \(y\), and use the binomial
coefficients \(1,5,10,10,5,1\).  The coefficient of \(y^4\) is
\(-B/A+B/A=0\); the remaining coefficients normalize to the four displayed
fractions.  The executable verifier checks the fully symbolic numerator.
\end{proof}

The formal derivative of~\eqref{eq:general-quintic} is a quartic.  Thus all
critical points are roots of a polynomial whose factorization identity is
already covered by the checked Ferrari module.  For real rational coefficients,
a Sturm chain or derivative-based interval subdivision is a bounded exact
algorithm for isolating every real root.  This gives an independent effective
meaning to the source page's elementary observation that a quintic's turning
points are controlled by a quartic.

\subsection{A transparent reduction to principal form}

A principal quintic has the form
\begin{equation}\label{eq:principal-form}
 Y^5+c_2Y^2+c_1Y+c_0=0;
\end{equation}
its quartic and cubic terms vanish.  We now prove the quadratic Tschirnhaus
step by power sums.

Let \(x_1,\ldots,x_5\) be the roots of a monic depressed quintic in a formal
splitting algebra, and set
\[
 S_m=\sum_{i=1}^5x_i^m.
\]
Then \(S_1=0\).  If the coefficient \(p\) of \(x^3\) in
\eqref{eq:depressed-quintic} is zero, the equation is already principal.  Assume
therefore \(p\ne0\).  Newton's identity gives \(S_2=-2p\ne0\).

Define
\begin{equation}\label{eq:quadratic-tsch}
 Y_i=x_i^2+\alpha x_i+\beta,
 \qquad \beta=-\frac{S_2}{5}.
\end{equation}
Then
\begin{equation}\label{eq:Ysum}
 \sum_iY_i=S_2+5\beta=0.
\end{equation}
A second expansion gives
\begin{align}
 \sum_iY_i^2
 &=S_4+2\alpha S_3+(\alpha^2+2\beta)S_2+5\beta^2\notag\\
 &=S_4+2\alpha S_3+\alpha^2S_2-\frac{S_2^2}{5}.
 \label{eq:Ysquare-sum}
\end{align}
Choose \(\alpha\) to satisfy
\begin{equation}\label{eq:alpha-principal}
 5S_2\alpha^2+10S_3\alpha+5S_4-S_2^2=0.
\end{equation}
This is a genuine quadratic because \(S_2\ne0\), and its required square root
is supplied by the checked quadratic factorization theorem.

\begin{theorem}[Principal-form reduction]\label{thm:principal-reduction}
With \(\alpha,\beta\) satisfying
\eqref{eq:quadratic-tsch}--\eqref{eq:alpha-principal}, the transformed roots
\(Y_i\) are the roots of a polynomial of the form~\eqref{eq:principal-form}.
The transformation and its correctness are expressible by finite radical-tower
certificates and are provable in \ipra.
\end{theorem}

\begin{proof}
Equations~\eqref{eq:Ysum} and~\eqref{eq:Ysquare-sum} show that the first two
power sums \(T_1=\sum Y_i\) and \(T_2=\sum Y_i^2\) vanish.  If
\(e_j(Y)\) denotes the elementary symmetric polynomial in the \(Y_i\), Newton's
identities give
\[
 e_1=T_1=0,
 \qquad
 2e_2=e_1T_1-T_2=0.
\]
Characteristic zero permits cancellation of \(2\).  Therefore
\[
 \prod_{i=1}^5(Z-Y_i)
 =Z^5-e_3(Y)Z^2+e_4(Y)Z-e_5(Y),
\]
which is principal.  Every symmetric coefficient can either be computed as a
bounded sum of products or obtained as a resultant.  The proof uses only the
quadratic relation~\eqref{eq:alpha-principal} and polynomial identities.
\end{proof}

The two possible roots of~\eqref{eq:alpha-principal} give two valid principal
forms.  This explains constructively the source page's statement that either
solution may be used.  No choice principle is needed: a certificate carries one
chosen radical and its defining square equation.

\subsection{Bring--Jerrard reduction as a compact certificate}

Let
\[
 P(Y)=Y^5+c_2Y^2+c_1Y+c_0
\]
a principal quintic.  A quartic Tschirnhaus map has the form
\begin{equation}\label{eq:quartic-tsch}
 T(Y)=Y^4+\alpha Y^3+\beta Y^2+\gamma Y+\delta.
\end{equation}
Introduce a new indeterminate \(V\) and compute
\begin{equation}\label{eq:tsch-resultant}
 R(V)=\Res_Y\bigl(P(Y),V-T(Y)\bigr).
\end{equation}
Because the first polynomial has degree five, \(R\) is monic of degree five in
\(V\), counting multiplicities.  A successful Bring certificate consists of
\[
 (\alpha,\beta,\gamma,\delta,d_1,d_0).
\]
together with their finite algebraic-tower definitions and the exact identity
\begin{equation}\label{eq:bring-certificate}
 R(V)=V^5+d_1V+d_0.
\end{equation}

\begin{theorem}[Soundness of a Bring--Jerrard certificate]
\label{thm:bring-certificate-sound}
If the finite identity~\eqref{eq:bring-certificate} holds and \(y\) is a root
of \(P\), then \(v=T(y)\) satisfies
\[
 v^5+d_1v+d_0=0.
\]
Checking the certificate is primitive recursive, uniformly in the coded
coefficient tower, and its soundness is provable in \ipra.
\end{theorem}

\begin{proof}
The defining property of the resultant is already visible from the Sylvester
determinant: substituting \(V=T(y)\) makes the two polynomials
\(P(Y)\) and \(T(y)-T(Y)\) have the common root \(Y=y\), so the determinant
vanishes.  Substitution into~\eqref{eq:bring-certificate} gives the desired
trinomial equation.  Alternatively a Bézout certificate
\(A(Y,V)P(Y)+B(Y,V)(V-T(Y))=R(V)\) makes this a single polynomial evaluation.
All matrices and degree bounds are finite.
\end{proof}

\begin{frontier}[Universal witness generation]
The classical Bring construction chooses the four parameters in
\eqref{eq:quartic-tsch} so that the coefficients of \(V^4,V^3,V^2\) in
\eqref{eq:tsch-resultant} vanish.  After suitable elimination, the remaining
choices require a cubic and a quadratic equation.  This is entirely compatible
with \ipra{} because the Cardano and quadratic factorization identities are
already checked.  What is large is not the logical principle but the universal
symbolic witness: expanded formulas can occupy megabytes.  This packet proves
the uniform verifier and soundness theorem, and gives an exact route for adding
any generated witness, but it does not claim a kernel-checked universal
multimegabyte certificate generator.
\end{frontier}

This distinction is important.  The mathematically useful theorem is often the
small verifier~\eqref{eq:bring-certificate}, not a human-readable expansion of
all its witnesses.

\subsection{Scaling a Bring--Jerrard trinomial}

The one-variable reduction is short and complete.

\begin{theorem}[Scaling formula]\label{thm:bring-scaling}
Let \(c^4=p\) and \(c\ne0\).  Then
\begin{equation}\label{eq:scaling-identity}
 (cz)^5+p(cz)+q
 =c^5\left(z^5+z+\frac{q}{c^5}\right).
\end{equation}
Consequently every compatible branch of a root of \(x^5+px+q\) has the form
\begin{equation}\label{eq:root-by-BR}
 x=c\,\BR\!\left(\frac{q}{c^5}\right),
\end{equation}
with the other four roots obtained by the other algebraic branches.  In the
customary notation this is
\(p^{1/4}\BR(p^{-5/4}q)\).  If \(p=0\), the equation is simply
\(x^5=-q\).
\end{theorem}

\begin{proof}
Use \(p=c^4\) in the left side and factor out \(c^5\).  Formula
\eqref{eq:root-by-BR} follows from the defining equation of the selected Bring
branch.  To avoid division inside \ipra{}, multiply by \(c^5\) and retain the
guard \(c\ne0\).
\end{proof}

Choosing \(c^4=-p\) instead produces the alternative form
\(z^5-z+a=0\).  The four choices of the fourth root and the five algebraic
branches explain why the displayed scaling formula must be accompanied by
branch compatibility; it is not a single global complex-valued radical.

\subsection{Inverting the Tschirnhaus transformation}

Suppose a root \(v\) of the Bring--Jerrard equation has been computed.  The
corresponding original root \(y\) is a common zero of
\begin{equation}\label{eq:inverse-system}
 P(Y)=0,
 \qquad T(Y)-v=0.
\end{equation}
The polynomial gcd of these two polynomials over the finite coefficient tower
can be computed by the subresultant algorithm.  In the generic case the gcd
has degree at most four, so its roots are available from the checked formulas
through Ferrari.  In a degenerate case the certificate records the actual gcd
and its degree; no generic division is silently performed.

\begin{proposition}[Certified inversion]
A list of candidates \(y_j\) is a complete inverse image of \(v\) under the
Tschirnhaus map on the roots of \(P\) when the following finite conditions
hold:
\begin{enumerate}
\item every \(y_j\) satisfies both equations in~\eqref{eq:inverse-system};
\item the monic gcd of the two left sides factors as \(\prod_j(Y-y_j)\);
\item the supplied Bézout/subresultant trace verifies that this polynomial is
      the gcd.
\end{enumerate}
These conditions have a primitive-recursive verifier and imply completeness in
\ipra.
\end{proposition}

This is the constructive resolution of the ``extraneous solutions'' issue
mentioned in the source page: numerical guessing is not logically necessary.
One may approximate candidates to choose a branch, but exact substitution and
the gcd certificate perform the final verification.

\subsection{Brioschi and other normal forms}

The Brioschi family
\[
 W^5-10CW^3+45C^2W-C^2=0.
\]
afforded a different route in the classical literature.  The source page gives
a rational Tschirnhaus map
\[
 W=\frac{\lambda+\mu X}{X^2/C-3}.
\]
For any proposed formulas for \(C,\lambda,\mu\), correctness is again a finite
certificate: multiply by the fifth power of the denominator, reduce modulo the
original quintic, and check that the remainder is zero.  The denominator guard
is retained, and exceptional roots are checked separately.  Thus the Brioschi
normal-form claim has the same compact \ipra{} interface as
Theorem~\ref{thm:bring-certificate-sound}, even though deriving the universal
polyhedral formulas for the parameters belongs to the global icosahedral
frontier of Section~\ref{sec:beyond}.

\section{Abel--Ruffini and finite Galois certificates}\label{sec:abel-ruffini}

The statement ``the quintic cannot be solved by radicals'' has two distinct
forms.  The generic theorem says that no single radical formula works for the
universal quintic.  A concrete theorem says that a specified irreducible
quintic has no radical expression for its roots.  Both can be organized
constructively, but neither says that every quintic is insoluble: reducible and
special solvable quintics are treated in Section~\ref{sec:solvable}.

\subsection{Radical towers yield solvable finite groups}

A finite radical certificate over a field \(K\) is a tower
\begin{equation}\label{eq:radical-tower}
 K=K_0\subset K_1\subset\cdots\subset K_m,
 \qquad K_{j+1}=K_j(\alpha_j),
 \qquad \alpha_j^{n_j}\in K_j,
\end{equation}
with the required roots of unity adjoined.  For object-theory purposes a finite
extension is represented by a basis, addition and multiplication tables,
embeddings of the preceding field, and reduction certificates for the defining
polynomials.  Every axiom is a bounded equality on table entries.

\begin{lemma}[Kummer step]\label{lem:kummer-step}
Suppose \(K\) contains all \(n\)-th roots of unity and
\(L=K(\alpha)\) with \(\alpha^n\in K\).  Every \(K\)-automorphism of the normal
closure sends \(\alpha\) to \(\zeta\alpha\) for an \(n\)-th root of unity
\(\zeta\).  The corresponding relative automorphism group embeds into the
abelian group of \(n\)-th roots of unity.
\end{lemma}

\begin{proof}
If \(\sigma\) fixes \(K\), then
\(\sigma(\alpha)^n=\sigma(\alpha^n)=\alpha^n\).  When \(\alpha\ne0\), the
quotient \(\sigma(\alpha)/\alpha\) is an \(n\)-th root of unity; the zero case
is trivial.  Multiplication of these quotients follows composition of
automorphisms.  In a coded finite extension, all divisions are implemented by
multiplicative-inverse tables with nonzero guards.
\end{proof}

\begin{theorem}[Radicals imply solvability]\label{thm:radicals-solvable}
The Galois group of the normal closure of a finite radical tower has a finite
normal series with abelian quotients.  Hence every polynomial solvable by
radicals has a solvable Galois group.
\end{theorem}

\begin{proof}
First adjoin the finitely many roots of unity required by the tower; these are
cyclotomic extensions with abelian Galois groups.  Normalize one radical step at
a time.  Lemma~\ref{lem:kummer-step} makes the kernel contributed by that step
abelian.  Induction on the tower length and the elementary fact that an
extension of a solvable group by a solvable group is solvable give the result.
For finite multiplication tables, normality, quotient multiplication, and the
derived series are all bounded checks.  The argument therefore translates into
\ipra{} for every coded tower.
\end{proof}

The converse---a polynomial with solvable Galois group is solvable by
radicals after adjoining roots of unity---uses the same series in reverse and
successive cyclic extensions.  In a fully explicit finite extension it is
again a certificate theorem.  We use both directions below for the standard
criterion ``solvable by radicals iff the finite Galois group is solvable.''

\subsection{The generic quintic has group \(S_5\)}

Let \(x_1,\ldots,x_5\) be algebraically independent indeterminates over
\(\Q\), let \(e_j\) be their elementary symmetric polynomials, and put
\[
 L=\Q(x_1,\ldots,x_5),
 \qquad K=\Q(e_1,\ldots,e_5).
\]
Every permutation of the \(x_i\) fixes \(K\), so \(S_5\) acts by
\(K\)-automorphisms of \(L\).  Conversely the roots of
\begin{equation}\label{eq:generic-quintic}
 Z^5-e_1Z^4+e_2Z^3-e_3Z^2+e_4Z-e_5
\end{equation}
are exactly the \(x_i\), so a \(K\)-automorphism permutes them.

The only subtle point is that the fixed field of the permutation action is
exactly \(K\).  This admits a finite algorithmic proof.

\begin{lemma}[Constructive symmetric rational functions]
\label{lem:symmetric-rational}
Every rational function in \(x_1,\ldots,x_5\) fixed by all permutations is a
rational function in \(e_1,\ldots,e_5\).  For bounded numerator and denominator
degrees, the conversion is primitive recursive.
\end{lemma}

\begin{proof}
Write the invariant as \(f=p/q\).  Replace the denominator by its permutation
norm
\[
 Q=\prod_{\sigma\in S_5}\sigma(q).
\]
Then \(Q\) is symmetric and \(P=fQ=p(Q/q)\) is a symmetric polynomial.
The fundamental theorem of symmetric polynomials has a terminating
lexicographic reduction: subtract from a symmetric polynomial a suitable
monomial in \(e_1,\ldots,e_5\) that cancels its leading monomial, and repeat.
The leading exponent strictly decreases in a finite well-order bounded by the
input degree.  Thus \(P=P(e_1,\ldots,e_5)\) and similarly for \(Q\).  The
reduction loop and its evaluation invariant are primitive recursive; the
termination measure is a bounded natural code of the exponent vector.
\end{proof}

\begin{theorem}[Generic Galois group]\label{thm:generic-S5}
The Galois group of the generic quintic~\eqref{eq:generic-quintic} over
\(K\) is \(S_5\).
\end{theorem}

\begin{proof}
The preceding discussion embeds the group in \(S_5\), while every permutation
of the indeterminates is a \(K\)-automorphism.  Lemma
\ref{lem:symmetric-rational} shows there is no larger fixed field hidden
between \(K\) and the invariant field, so the action is the full Galois action.
All coefficient manipulations at a requested degree are finite.  This is the
constructive core of the generic argument developed in ProVeIt's
polynomial-formulas/Abel--Ruffini packet.
\end{proof}

\subsection{Why \(S_5\) is not solvable}

Let the commutator convention be
\([g,h]=ghg^{-1}h^{-1}\).

\begin{lemma}\label{lem:S5-derived}
The commutator subgroup of \(S_5\) is \(A_5\).
\end{lemma}

\begin{proof}
The sign homomorphism has abelian target, so every commutator is even and
\([S_5,S_5]\subseteq A_5\).  On the other hand, the commutator of two
transpositions sharing one letter is a 3-cycle (up to orientation), and
3-cycles generate \(A_5\).  Hence \(A_5\subseteq[S_5,S_5]\).
\end{proof}

\begin{lemma}\label{lem:A5-perfect}
The group \(A_5\) is perfect: \([A_5,A_5]=A_5\).
\end{lemma}

\begin{proof}
Both
\[
 g=(2\,3)(4\,5),
 \qquad h=(1\,2)(4\,5)
\]
are in \(A_5\), and a direct five-entry permutation calculation gives
\[
 [g,h]=(1\,2\,3).
\]
The commutator subgroup is normal.  All 3-cycles are conjugate within
\(A_5\), and 3-cycles generate \(A_5\), so the commutator subgroup is all of
\(A_5\).
\end{proof}

Both lemmas are tiny finite certificates: the archive independently enumerates
the 120 permutations and obtains derived-series orders
\[
 120,60,60.
\]

\begin{theorem}[Constructive Abel--Ruffini theorem]
\label{thm:abel-ruffini}
There is no radical expression, uniform in the five generic coefficients, that
produces a root of the generic quintic.
\end{theorem}

\begin{proof}
By Theorem~\ref{thm:generic-S5}, the generic Galois group is \(S_5\).
Lemmas~\ref{lem:S5-derived} and~\ref{lem:A5-perfect} show that its derived
series never reaches the trivial group, so it is not solvable.  This contradicts
Theorem~\ref{thm:radicals-solvable} for any proposed radical tower.  Every group
calculation is finite, and the symmetric-polynomial reduction is a bounded
algorithm, so the argument has the required \ipra{} form.
\end{proof}

The theorem excludes a universal formula by ordinary radicals.  It does not
exclude Bring radicals, hypergeometric functions, numerical algorithms, or
radical formulas for special quintics.

\subsection{A finite group criterion for an irreducible quintic}

An irreducible quintic over \(\Q\) has a transitive Galois group
\(G\le S_5\).  The solvable transitive subgroups have a particularly small
form.

\begin{theorem}[Prime-degree solvable subgroup criterion]
\label{thm:AGL-criterion}
A solvable transitive subgroup of \(S_5\) is conjugate to a subgroup of
\begin{equation}\label{eq:AGL}
 \operatorname{AGL}(1,5)
 =\{x\mapsto ax+b:a\in\F_5^\times,\ b\in\F_5\},
\end{equation}
which has order \(20\).  Conversely every subgroup of
\(\operatorname{AGL}(1,5)\) is solvable.  Hence an irreducible quintic is
solvable by radicals exactly when its Galois group is conjugate to such a
subgroup.
\end{theorem}

\begin{proof}
Let \(G\) be solvable and transitive, and choose a nontrivial minimal normal
subgroup \(N\).  A standard finite induction on the derived series shows that
\(N\) is elementary abelian of prime-power order.  The orbits of a normal
subgroup form a block system.  Since the action has prime degree five, a
nontrivial normal subgroup must be transitive.  Therefore \(5\mid |N|\).
Inside \(S_5\), a 5-subgroup has order at most five, so \(N\cong C_5\) and acts
regularly.  Conjugation embeds a point stabilizer into
\(\Aut(C_5)\cong C_4\), giving
\(G\le C_5\rtimes C_4=\operatorname{AGL}(1,5)\) after relabeling the points.

Conversely the translation subgroup \(C_5\) is normal in
\(\operatorname{AGL}(1,5)\) and the quotient is \(C_4\), so the affine group
and all its subgroups are solvable.  The radical criterion follows from
Theorem~\ref{thm:radicals-solvable} and its converse.
\end{proof}

The executable verifier checks the theorem independently by enumerating all
156 subgroups of \(S_5\): there are 20 transitive subgroups as actual subsets,
18 solvable ones of orders \(5,10,20\), and two nonsolvable ones, \(A_5\) and
\(S_5\).  The affine group's derived-series orders are \(20,5,1\).

\subsection{A concrete nonsolvable quintic: \(x^5-x+1\)}

Let
\[
 f(X)=X^5-X+1.
\]
Its discriminant, from~\eqref{eq:trinomial-discriminant}, is
\[
 \Disc(f)=-256+3125=2869.
\]
Modulo two one has the exact square-free factorization
\begin{equation}\label{eq:f-mod2}
 f(X)=(X^2+X+1)(X^3+X^2+1)\pmod2.
\end{equation}
The quadratic and cubic factors have no root in \(\F_2\), hence are
irreducible.  The factor degrees are therefore \((2,3)\).

Modulo three, define \(r_0=X\) and
\(r_{j+1}=r_j^3\bmod f\).  A finite reduction gives
\begin{center}
\begin{tabular}{c|l}
\toprule
\(j\)&\(r_j\) in \(\F_3[X]/(f)\)\\
\midrule
1&\(X^3\)\\
2&\(-X^4+X-1\)\\
3&\(-X^4-X^2-1\)\\
4&\(-X^4-X^3+X^2+X-1\)\\
5&\(X\)\\
\bottomrule
\end{tabular}
\end{center}
Moreover the Bézout identity
\begin{equation}\label{eq:f-bezout3}
 f(X)+(-X^2-1)(X^3-X)=1\pmod3
\end{equation}
shows that \(f\) has no linear factor over \(\F_3\).  Since
\(f\mid X^{3^5}-X\) by the table, and five is prime, all irreducible factors
have degree five.  Thus \(f\) is irreducible modulo three and hence over
\(\Q\).

\begin{theorem}\label{thm:x5-x+1-S5}
The Galois group of \(X^5-X+1\) over \(\Q\) is \(S_5\).  In particular the
polynomial is not solvable by radicals.
\end{theorem}

\begin{proof}
The factorization modulo three supplies a Frobenius element of cycle type
\((5)\), while~\eqref{eq:f-mod2} supplies one of type \((2)(3)\).  Cubing the
latter leaves a transposition.  A subgroup of \(S_5\) containing a 5-cycle and
a transposition is all of \(S_5\): conjugating the transposition by powers of
the 5-cycle gives the edge transpositions of a connected graph on the five
letters, and edge transpositions of a connected graph generate the full
symmetric group.  Nonsolvability follows from
Lemmas~\ref{lem:S5-derived}--\ref{lem:A5-perfect}.
\end{proof}

The factor-degree/Frobenius lemma used here has a finite \ipra{} realization.
For a prime not dividing the discriminant, reduce the integral splitting order
modulo the prime, choose a prime ideal above it by bounded search in the finite
quotient, and use the residue Frobenius \(z\mapsto z^p\).  Its orbits on the
roots have lengths equal to the irreducible factor degrees.  Addition and
multiplication in every residue algebra are finite tables; the decomposition
and lifting checks are therefore certificate verifications.  This is a
substantial Lean module, but it needs no infinitary set existence.

\subsection{The real number \(\BR(3)\) is not a radical number}

Now let
\[
 g(X)=X^5+X+3.
\]
The discriminant is
\[
 \Disc(g)=256+3125\cdot3^4=253381.
\]
Modulo two, \(g\) has the same factorization~\eqref{eq:f-mod2}.  Modulo seven,
let \(s_0=X\) and \(s_{j+1}=s_j^7\bmod g\).  Exact reduction gives
\begin{center}
\begin{tabular}{c|l}
\toprule
\(j\)&\(s_j\) in \(\F_7[X]/(g)\)\\
\midrule
1&\(-X^3-3X^2\)\\
2&\(3X^4+2X^3+3X+1\)\\
3&\(2X^4-3X^3+3X^2+2X+3\)\\
4&\(2X^4+2X^3+X+3\)\\
5&\(X\)\\
\bottomrule
\end{tabular}
\end{center}
Writing \(h=s_1-X=-X^3-3X^2-X\), the explicit identity
\begin{equation}\label{eq:g-bezout7}
 -2g(X)+(-2X^2-X-2)h(X)=1\pmod7
\end{equation}
excludes linear factors.  As above, \(g\) is irreducible modulo seven.

\begin{theorem}\label{thm:BR3-not-radical}
The Galois group of \(X^5+X+3\) is \(S_5\).  Its unique real root
\(\BR(3)\) is not contained in any radical extension of \(\Q\).
\end{theorem}

\begin{proof}
Modulo seven the factor degree is five, giving a 5-cycle; modulo two the factor
degrees are two and three, giving a transposition after cubing.  Thus the group
is \(S_5\).  If the single algebraic number \(\BR(3)\) belonged to a radical
tower, the normal closure of that tower (after adjoining roots of unity) would
have solvable Galois group.  The splitting field of the minimal polynomial of
\(\BR(3)\) is a subextension of that normal closure, so its Galois group would
be a quotient of a solvable group and hence solvable, contradicting \(S_5\).
\end{proof}

This gives a concrete arithmetic interpretation of the source page's claim
that the Bring radical is not an ordinary radical function.  The generic
version follows already from Theorem~\ref{thm:abel-ruffini}; the specialization
at \(a=3\) shows that the obstruction is not merely formal.

\section{Solvable quintics and cyclotomic resolvents}\label{sec:solvable}

Abel--Ruffini leaves a large and interesting solvable class.  In this section
all displayed root identities are proved by finite expansion in a cyclotomic
extension.  The point is not that the resulting formulas are small, but that
their correctness is decidable by exact algebra.

\subsection{Reducible equations}

A reducible quintic is solved factor by factor.  For example,
\begin{equation}\label{eq:reducible-example}
 X^5-X^4-X+1=(X^2+1)(X+1)(X-1)^2.
\end{equation}
This identity is immediate by multiplication and is an \ipra{} theorem after
integer polynomial normalization.  Its roots use only the checked quadratic
formula and rational roots.  Such examples are not exceptions to
Theorem~\ref{thm:abel-ruffini}: that theorem concerns a uniform formula for the
generic quintic.

\subsection{The de Moivre quintic}

Let \(\omega\) satisfy~\eqref{eq:cyclo5}, and let \(u,v\) lie in a commutative
field containing \(\omega\).  Define
\[
 r_j=\omega^j u+\omega^{-j}v,
 \qquad 0\le j<5.
\]

\begin{theorem}[de Moivre product identity]\label{thm:de-moivre-product}
One has
\begin{equation}\label{eq:de-moivre-product}
 \prod_{j=0}^{4}(X-r_j)
 =X^5-5uvX^3+5u^2v^2X-(u^5+v^5).
\end{equation}
This identity is valid in the finite quotient
\(\Z[\omega,u,v,X]/(1+\omega+\cdots+\omega^4)\) and is therefore directly
formalizable in \ipra.
\end{theorem}

\begin{proof}
Let \(s_m=\sum_jr_j^m\).  Since
\(\sum_{j=0}^4\omega^{kj}=0\) unless \(5\mid k\), expansion gives
\[
 s_1=s_3=0,
 \qquad s_2=10uv,
 \qquad s_4=30u^2v^2.
\]
Newton's identities then yield
\[
 e_1=0,
 \quad e_2=-5uv,
 \quad e_3=0,
 \quad e_4=5u^2v^2.
\]
Finally,
\[
 e_5=\prod_jr_j
 =\prod_j(v+u\omega^{2j})=u^5+v^5,
\]
because multiplication by two permutes the fifth roots of unity and
\(\prod_{j=0}^4(Z+\omega^j)=Z^5+1\).  Substitution into
\(\prod(X-r_j)=X^5-e_1X^4+e_2X^3-e_3X^2+e_4X-e_5\) gives the result.
\end{proof}

Take \(uv=-a\) and \(u^5+v^5=-b\).  Then~\eqref{eq:de-moivre-product} becomes
\begin{equation}\label{eq:de-moivre-quintic}
 X^5+5aX^3+5a^2X+b.
\end{equation}
If \(Y=u^5\), then \(v=-a/u\) and
\[
 v^5=-\frac{a^5}{Y}.
\]
The equation \(u^5+v^5=-b\) is therefore equivalent to
\begin{equation}\label{eq:de-moivre-aux}
 Y^2+bY-a^5=0.
\end{equation}

\begin{corollary}[Radical solution of the de Moivre family]
\label{cor:de-moivre-roots}
Let \(Y\) be either root of~\eqref{eq:de-moivre-aux}, choose \(u^5=Y\), and
put \(v=-a/u\).  Then the five roots of~\eqref{eq:de-moivre-quintic} are
\begin{equation}\label{eq:de-moivre-roots}
 X_j=\omega^ju-\frac{a}{\omega^ju},
 \qquad 0\le j<5.
\end{equation}
Every defining relation and the final substitution are finite radical-tower
certificates.
\end{corollary}

The auxiliary equation is quadratic; the remaining operations are fifth roots
and fifth roots of unity.  Thus this family is solved by ordinary radicals in
the strict sense.

\subsection{The four-resolvent identity printed by MathWorld}

Let \(u_1,u_2,u_3,u_4\) be elements of a commutative cyclotomic extension and
put
\begin{equation}\label{eq:four-resolvent-roots}
 r_j=\omega^ju_1+\omega^{2j}u_2+\omega^{3j}u_3+\omega^{4j}u_4.
\end{equation}
Define
\begin{align*}
 U&=u_1u_4+u_2u_3,\\
 V&=u_1u_2^2+u_2u_4^2+u_3u_1^2+u_4u_3^2,\\
 W&=u_1^2u_4^2+u_2^2u_3^2-u_1^3u_2-u_2^3u_4-u_3^3u_1-u_4^3u_3
     -u_1u_2u_3u_4,\\
 X_0&=u_1^3u_3u_4+u_2^3u_1u_3+u_3^3u_2u_4+u_4^3u_1u_2,\\
 Y_0&=u_1u_3^2u_4^2+u_2u_1^2u_3^2+u_3u_2^2u_4^2+u_4u_1^2u_2^2,\\
 Z&=u_1^5+u_2^5+u_3^5+u_4^5.
\end{align*}

\begin{theorem}[Four-resolvent product]\label{thm:four-resolvent}
The five values~\eqref{eq:four-resolvent-roots} are the roots of
\begin{equation}\label{eq:four-resolvent-product}
 T^5-5UT^3-5VT^2+5WT+5(X_0-Y_0)-Z.
\end{equation}
\end{theorem}

\begin{proof}
Expand \(\prod_{j=0}^4(T-r_j)\), reduce every power of \(\omega\) with
\(1+\omega+\cdots+\omega^4=0\), and collect monomials in the \(u_i\).  The
coefficients are exactly those displayed above.  This finite calculation is
implemented independently by the archive's checker; its soundness in \ipra{}
is Lemma~\ref{lem:poly-normalize} in the cyclotomic quotient.  A human
alternative is to compute the first five power sums of the \(r_j\) and apply
Newton's identities, using the root-of-unity filter as in
Theorem~\ref{thm:de-moivre-product}.
\end{proof}

\begin{corollary}[Conditional root formula for \(T^5+aT+b\)]
\label{cor:four-resolvent-specialization}
If
\begin{equation}\label{eq:four-resolvent-constraints}
 U=0,
 \qquad V=0,
 \qquad W=\frac a5,
 \qquad 5(X_0-Y_0)-Z=b,
\end{equation}
then every \(r_j\) in~\eqref{eq:four-resolvent-roots} satisfies
\(r_j^5+ar_j+b=0\).
\end{corollary}

This proves the algebraic core of MathWorld's formulas (43)--(55).  The long
expressions there for the \(u_i\) are a way of manufacturing a solution of
\eqref{eq:four-resolvent-constraints}.  Their correctness can be checked by
adjoining symbols for the indicated square and fifth roots, clearing
nonzero denominators, and normalizing the four identities.  That is a finite
certificate problem; the present packet proves the uniform product identity
and specialization, while leaving the page's full parameter substitution as a
separate generated certificate rather than reproducing several pages of
radicals.

\subsection{Why every solvable irreducible quintic has resolvents of this kind}

The affine-group criterion gives a conceptual construction.  Label the roots
\(x_0,\ldots,x_4\) so that the normal translation subgroup \(C_5\) acts by
\(x_j\mapsto x_{j+1}\).  Define Lagrange resolvents
\begin{equation}\label{eq:lagrange-resolvents}
 R_k=\sum_{j=0}^4\omega^{-kj}x_j,
 \qquad 0\le k<5.
\end{equation}
Under the generating 5-cycle,
\[
 R_k\longmapsto\omega^kR_k.
\]
Therefore \(R_k^5\) is fixed by the normal \(C_5\).  The quotient of the
Galois group by \(C_5\) is a subgroup of \(C_4\), so each \(R_k^5\) lies in an
extension of degree at most four over the base field and can be obtained using
the checked formulas through Ferrari.  Once compatible fifth roots \(R_k\)
are chosen, Fourier inversion gives
\begin{equation}\label{eq:resolvent-inversion}
 x_j=\frac15\sum_{k=0}^4\omega^{kj}R_k.
\end{equation}

\begin{theorem}[Constructive resolvent solution]
\label{thm:solvable-quintic-resolvents}
Given an irreducible quintic together with a finite Galois certificate embedding
its group into \(\operatorname{AGL}(1,5)\), the construction
\eqref{eq:lagrange-resolvents}--\eqref{eq:resolvent-inversion} yields a finite
radical-tower certificate for all five roots.  Every compatibility choice is
verified by substitution rather than by an unbounded search.
\end{theorem}

\begin{proof}
The normal \(C_5\) and quotient action are read from the supplied permutation
table.  Elementary symmetric polynomials in the \(R_k^5\) are fixed by the
whole group and can be computed as rational functions of the original
coefficients by finite orbit sums.  The quotient degree is at most four, so the
checked quadratic, Cardano and Ferrari identities factor the corresponding
resolvent polynomial.  Adjoin one fifth root and obtain the remaining
compatible roots by the recorded group action.  Finally verify
\eqref{eq:resolvent-inversion} and the original quintic by cyclotomic
normalization.  Each stage is a bounded calculation on finite codes.
\end{proof}

This theorem explains the source page's statement that for prime degree five a
root may be written as a sum of fifth roots of the roots of a quartic auxiliary
equation, with compatibility relations among those fifth roots.

\subsection{Explicit radical examples}

The following examples illustrate how short the final verification can be even
when discovering the expression was difficult.

\begin{example}[A single fifth radical]
Let \(t^5=2\) and
\[
 x=1+t-t^2+t^3-t^4.
\]
Then
\begin{equation}\label{eq:explicit-example-one}
 x^5-5x^4+30x^3-50x^2+55x-21=0.
\end{equation}
Indeed, substituting the displayed \(x\) makes the left side a multiple of
\(t^5-2\).  Polynomial division gives an exact quotient of degree 15, recorded
by the executable certificate.
\end{example}

\begin{example}[A Paxton Young example]
With the same \(t^5=2\), put
\[
 x=t-t^2+t^3-t^4.
\]
Then
\begin{equation}\label{eq:young-example}
 x^5+20x^3+20x^2+30x+10=0.
\end{equation}
Again the substituted polynomial has zero remainder modulo \(t^5-2\).
\end{example}

\begin{example}[An eleventh-root-of-unity quintic]
Let \(z^{11}=1\), \(z\ne1\), and \(x=z+z^{-1}\).  Then
\begin{equation}\label{eq:cos11-example}
 x^5+x^4-4x^3-3x^2+3x+1=0.
\end{equation}
To verify this, multiply the left side after substituting \(x=z+z^{-1}\) by
\(z^5\) and reduce modulo
\(1+z+z^2+\cdots+z^{10}\).  The remainder is zero.  Under the usual complex
embedding, these roots are \(2\cos(2k\pi/11)\).
\end{example}

All three examples are exact quotient-ring computations and require no appeal
to numerical approximation.

\subsection{Fifth roots of unity by square roots}

The de Moivre formulas themselves use \(\omega\), but a primitive fifth root of
unity is obtainable by quadratic radicals.  Let \(s^2=5\), choose
\(\alpha,\beta\in\{-1,1\}\), let
\(t^2=-10-2\beta s\), and put
\begin{equation}\label{eq:fifth-root-radical}
 \omega_{\alpha,\beta}
 =\frac{\alpha t+\beta s-1}{4}.
\end{equation}
Clearing the denominator and reducing successively by
\(s^2-5\) and \(t^2+10+2\beta s\) proves
\[
 1+\omega_{\alpha,\beta}+\omega_{\alpha,\beta}^2
 +\omega_{\alpha,\beta}^3+\omega_{\alpha,\beta}^4=0.
\]
The four sign pairs give the four primitive roots; inequality of their real and
imaginary components is decidable in the corresponding ordered quadratic
subfields.  Thus every cyclotomic constant used above is itself represented by
a short radical certificate.

\subsection{Parameterized solvable trinomials}

The source pages quote the Runge--Cadenhad--Young and Spearman--Williams
parameterizations of solvable irreducible trinomials \(X^5+aX+b\).  Their
logical content separates naturally:

\begin{enumerate}
\item The displayed rational formulas for \(a,b\) are primitive-recursive
      maps of the rational parameters.
\item The displayed radical expressions for the \(u_i\), once supplied, are
      checked by Corollary~\ref{cor:four-resolvent-specialization} after
      clearing their square- and fifth-root relations.
\item The ``if and only if'' classification is a Galois theorem.  By
      Theorem~\ref{thm:AGL-criterion}, it reduces to constructing rational
      invariants that parameterize the embeddings of the particular trinomial
      Galois group into \(\operatorname{AGL}(1,5)\).  This is finite invariant
      theory, but its published formulas are lengthy.
\end{enumerate}

Accordingly, this packet proves the finite group criterion and the universal
root verifier, and treats the exact necessity/sufficiency substitution for the
published parameterization as a generated certificate target in the Lean
roadmap.  This is stronger than accepting the formulas numerically, but more
honest than claiming a handwritten formal derivation that has not been
expanded.

\subsection{Casus irreducibilis}

The source page notes that some solvable quintics have five real roots but any
radical formula passes through nonreal intermediate values.  The field-theoretic
mechanism is the same as for the cubic.  If an irreducible, totally real
extension were generated through a tower in which every radical and every
intermediate embedding remained real, then any odd-prime radical step would
have only one real conjugate and could not contribute to the normal totally
real closure.  After eliminating such steps, only quadratic extensions remain,
so the degree is a power of two.  An irreducible quintic has degree five, not a
power of two; hence a radical solution of a five-real-root solvable quintic must
use complex intermediates.

For \ipra{} this argument is formulated for coded finite ordered fields and
embedding tables: at an odd-prime radical step the other roots are
\(\zeta^j\alpha\), and the table certifies that they are nonreal; a tower that
keeps every conjugate real can therefore have only index two steps.  The full
ordered-field tower module is included as a formalization target rather than a
completed Lean file, so the audit labels this global theorem \(\Outline\), even
though every individual finite tower can already be checked by the certificate
interface.

\section{Beyond radicals: what can be made constructive}\label{sec:beyond}

The reference pages devote substantial space to Hermite's theta-function
solution, Klein's icosahedral method, hypergeometric differential equations,
and the Doyle--McMullen iteration.  These claims are not excluded by
Abel--Ruffini: they use operations more general than ordinary radicals.  This
section separates their finite local cores from the global analytic work still
required for a complete Lean development.

\subsection{Named hypergeometric functions}

For fixed rational parameters \(a_1,\ldots,a_p,b_1,\ldots,b_q\), define
\[
 c_0=1,
 \qquad
 c_{n+1}=c_n
 \frac{(n+a_1)\cdots(n+a_p)}{(n+b_1)\cdots(n+b_q)(n+1)}.
\]
After clearing the fixed parameter denominators, this is a primitive-recursive
rational recurrence.  The truncated sum
\[
 H_N(z)=\sum_{n<N}c_nz^n
\]
is therefore a finite coded polynomial.  On a named disk where the coefficient
ratio is eventually bounded by a rational \(\theta<1\), the geometric tail
\[
 \left|H(z)-H_N(z)\right|
 \le \frac{|c_Nz^N|}{1-\theta}
\]
gives a primitive-recursive modulus.

\begin{proposition}[Local hypergeometric verification]
For fixed rational parameters, every coefficient identity, linear differential
equation, and algebraic substitution involving a generalized hypergeometric
series can be proved in \ipra{} coefficient by coefficient.  Analytic equality
on a named disk follows when both sides are supplied with effective tail
moduli.
\end{proposition}

\begin{proof}
The coefficient recurrence proves the differential equation after multiplying
by a common denominator and comparing the coefficient of \(z^n\).  Equality of
two named series follows by induction on \(n\) from equal initial coefficients
and a common recurrence.  A geometric or ratio-based tail modulus lifts the
finite equality to the analytic names.
\end{proof}

Theorem~\ref{thm:bring-hypergeom} is a complete example.  The four
hypergeometric basis functions used for all branches of the Bring quintic are
handled the same way: verify the Cockle--Harley recurrence, match four initial
data, and use the zero sum of the five roots.  This is a substantially smaller
task than formalizing a general library of analytic continuation.

The MathWorld observation that equations of the shape
\(x^p+bx^q+c=0\) admit one-variable hypergeometric expansions also has a
finite local interpretation.  Around a simple root or a parameter value where
the equation becomes \(x=X\Phi(x)\), finite Lagrange inversion gives the
coefficient sequence; residue classes of the exponent produce a finite number
of hypergeometric recurrences.  Each fixed pair \((p,q)\) is an \ipra{} theorem
schema.  Uniformity over arbitrary coded \((p,q)\) requires only bounded
factorial and binomial algorithms, not higher-order quantification.

\subsection{Theta series and elliptic data}

For \(|q|<1\), the Jacobi theta constants may be named by
\begin{align*}
 \theta_2(0,q)&=2q^{1/4}\sum_{n\ge0}q^{n(n+1)},\\
 \theta_3(0,q)&=1+2\sum_{n\ge1}q^{n^2}.
\end{align*}
On a named closed disk \(|q|\le r<1\), the omitted tail after \(N\) is bounded
by a geometric series beginning at \(r^{N^2}\).  Thus both theta constants are
primitive-recursively named.  If a rational lower bound for
\(|\theta_3(0,q)|\) is supplied, the modular lambda value
\begin{equation}\label{eq:lambda-theta}
 m(q)=\frac{\theta_2(0,q)^4}{\theta_3(0,q)^4}
\end{equation}
has an effective quotient name.

The complete elliptic integral
\[
 K(k)=\frac\pi2\,{}_2F_1\!\left(\frac12,\frac12;1;k^2\right).
\]
automatically receives a name from the preceding hypergeometric construction
on \(|k|<1\).  The nome
\[
 q=\exp\!\left(-\pi\frac{K(\sqrt{1-k^2})}{K(k)}\right).
\]
gains an effective name when the square roots, denominator lower bound, and
exponential modulus are supplied.  All of these operations are already within
the general effective-analysis architecture used elsewhere in ProveIt.

\begin{theorem}[Finite verification interface for Hermite-type formulas]
\label{thm:hermite-interface}
Fix a compact parameter region avoiding every displayed denominator and branch
cut, and supply:
\begin{enumerate}
\item regular names and moduli for the required theta constants, elliptic
      integrals, roots and exponentials;
\item the finite modular-equation identities relating the auxiliary modulus to
      the Bring parameter;
\item a branch certificate at one base point.
\end{enumerate}
Then the Hermite--Kronecker--Brioschi root expression can be verified in
\ipra{} by approximating both sides and checking the exact algebraic identity
\(x^5-x+a=0\), with uniqueness/continuation on the certified region.
\end{theorem}

\begin{proof}
Every special-function value has a primitive-recursive Cauchy name by the tail
bounds above.  The algebraic combination of names gives a candidate \(x\).
The modular equations and theta identities reduce the residual
\(x^5-x+a\) to zero; at finite precision this is an equality of truncated
series plus an explicit tail bound.  A lower bound for \(|5x^4-1|\) gives local
uniqueness by the quantitative inverse theorem.  Subdivision of a certified
path transports the branch from the base point to the target.
\end{proof}

\begin{frontier}[Global Hermite formula]
A complete proof of the exact formula printed on the source page must formalize
the relevant modular equations, all eighth-root branch choices, and analytic
continuation across a specified cut plane.  The local theta and hypergeometric
objects are fully compatible with \ipra{}, but this archive does not claim that
those global modular identities have already been reduced to Lean certificates.
The status is therefore \(\Outline\), not a rejection of the claim.
\end{frontier}

\subsection{Klein's icosahedral equation}

Klein reduces a principal quintic to an icosahedral invariant equation and
expresses a local inverse by a ratio of two Gauss hypergeometric functions.  The
large polynomial printed by MathWorld is itself an exact finite object.  For a
proposed value of its invariant \(Z\), substitution and polynomial
normalization are immediate \ipra{} tasks.

The hypergeometric ratio has three constructive layers.

\begin{enumerate}
\item Its numerator and denominator are named by rational coefficient
      recurrences, with effective tails for \(|1728Z|<1\).
\item A Wronskian or direct coefficient calculation verifies the associated
      second-order differential equation.
\item The rational icosahedral invariant of the ratio is checked by finite
      series substitution near a base point.  Equality then propagates on a
      certified simply connected region by uniqueness of the differential
      equation.
\end{enumerate}

The local disk is thus accessible with existing tools.  The global theorem that
all relevant branches and icosahedral monodromy are covered requires an
explicit path-and-monodromy library.  Such a library can be entirely
arithmetized: a path is a rational polygon, continuation uses overlapping
rational disks, and a permutation of the five roots is a five-entry table.
What remains is volume, not a need for classical set existence.

\subsection{Differential resolvents in arbitrary fixed degree}

The derivation of Theorem~\ref{thm:differential-resolvent} generalizes.  If
\(P(x,a)=0\) and \(P_x\ne0\), then
\[
 x'=-\frac{P_a}{P_x}.
\]
Repeated total differentiation expresses each \(x^{(j)}\) as a rational
function of \(x,a\).  The vector space over \(\Q(a)\) spanned by
\(1,x,\ldots,x^{d-1}\) has dimension \(d\) modulo a degree-\(d\) equation.
Therefore \(d+1\) derivatives are linearly dependent after reduction modulo
\(P\), producing a linear differential resolvent of order at most \(d\), and
o larger than \(d-1\) when the sum of the branches supplies one dependence.

For fixed \(d\), Gaussian elimination over rational-function codes and
polynomial remainder are primitive recursive.  Thus the source page's claim
that differential resolvents extend to higher-degree algebraic equations has a
uniform finite proof for each fixed degree.  Partial differential resolvents
for several coefficient variables are handled by multi-indexed bounded
derivatives in the same way.

\subsection{Doyle--McMullen iteration}

The source page describes a rational map \(T_Z\) for the Brioschi quintic and
advises iterating from a random start until convergence.  Randomness and the
phrase ``until it converges'' are not a formal theorem.  The correct
constructive interface is a trapping/contraction certificate.

\begin{theorem}[Certified rational iteration]\label{thm:certified-iteration}
Let \(T\) be a rational map with named coefficients.  Suppose a certificate
provides a closed rational disk \(D\), a nonzero lower bound for every
denominator on \(D\), a proof that \(T(D)\subseteq D\), and a rational
\(0\le\lambda<1\) such that
\[
 |T(u)-T(v)|\le\lambda|u-v|
 \quad(u,v\in D).
\]
Then iteration from any named point in \(D\) converges to the unique fixed point
with the primitive-recursive error bound
\[
 |T^n(w)-w_*|
 \le\frac{\lambda^n}{1-\lambda}|T(w)-w|.
\]
If the final algebraic reconstruction formulas are supplied, the resulting
Brioschi roots are effective \ipra{} reals or complexes.
\end{theorem}

\begin{proof}
The usual telescoping proof is a finite induction:
\(|w_{n+1}-w_n|\le\lambda^n|w_1-w_0|\).  A finite geometric sum bounds every
Cauchy tail, and a primitive-recursive logarithmic cutoff makes it smaller than
\(2^{-k}\).  Substitution into the rational map and then the Brioschi
polynomial verifies the limit.
\end{proof}

A global theorem that a generic ``random'' start reaches such a disk is a
separate dynamical statement involving Julia sets and measure.  It is not
needed for certified computation of any fixed instance and is not claimed in
this packet.

\section{Certified computation and the celestial-mechanics quintic}\label{sec:applications}

\subsection{At least one real root of every real quintic}

Let
\[
 P(x)=a_5x^5+a_4x^4+\cdots+a_0,
 \qquad a_5\ne0.
\]
a polynomial with rational coefficients.  Put
\[
 M=2+\sum_{j<5}\left|\frac{a_j}{a_5}\right|.
\]
For \(x=M\), the leading term dominates the sum of lower terms, and for
\(x=-M\) it has the opposite sign.  More explicitly,
\[
 \left|\sum_{j<5}a_jx^j\right|
 \le |a_5|(M-2)\sum_{j<5}M^j
 <|a_5|M^5.
\]
after replacing the displayed coarse bound by its elementary finite geometric
estimate.  Thus \(P(-M)P(M)<0\).  Exact rational bisection produces a regular
real root with a primitive-recursive modulus.

This proof is the constructive odd-degree theorem needed for quintics.  It does
not invoke the full fundamental theorem of algebra.  To find all real roots,
one may use either of two finite methods:

\begin{enumerate}
\item compute the quartic derivative, isolate its real roots using the checked
      lower-degree factorization identities, and bisect on the resulting
      monotonicity intervals; or
\item compute a signed subresultant/Sturm chain and use its exact variation
      count on rational endpoints.
\end{enumerate}

Both algorithms are primitive recursive for fixed degree five.  Multiple roots
are detected by the polynomial gcd \(\gcd(P,P')\).

\subsection{Certified evaluation of the Bring radical}

For the real Bring radical, bisection has two useful a posteriori certificates.
If a rational interval \([L,U]\) satisfies
\[
 P_a(L)\le0\le P_a(U).
\]
the root lies in the interval and its midpoint error is at most
\((U-L)/2\).  If a single rational candidate \(q\) is available, the sharper
residual estimate~\eqref{eq:residual-bound} gives
\[
 |q-\BR(a)|\le|q^5+q+a|.
\]
Thus a numerical library may use Newton, Halley, or a hypergeometric series for
speed and then certify the answer by one exact rational residual computation.
The executable artifact illustrates this with 80 exact bisection steps for
\(a=3\); it outputs a rational interval of width \(8/2^{80}\).

Inside \(|a|<R\), Theorem~\ref{thm:bring-radius} supplies a series tail bound.
Outside that disk the real branch remains easy: bisection is global because of
strong monotonicity.  Complex branches can be continued by overlapping
rational disks, each equipped with a lower bound for \(|5x^4+1|\) and a
quantitative inverse radius.  The branch history is a finite path certificate.

\subsection{Clearing the Lagrange-point denominators}

Let \(M_S,M_E>0\) be the two primary masses, let \(R>0\) be their separation,
and let \(r>0\) denote the distance of a collinear test mass from the second
primary.  Introduce \(\varepsilon\in\{+1,-1\}\) to encode the two sign
conventions.  The force-balance equation printed on the source page is
\begin{equation}\label{eq:lagrange-force}
 \frac{GM_S}{(R+\varepsilon r)^2}
 +\varepsilon\frac{GM_E}{r^2}
 =\omega^2(R+\varepsilon r).
\end{equation}
The satellite mass has already been cancelled.  Under the nonzero-denominator
guards \(r\ne0\) and \(R+\varepsilon r\ne0\), multiplication by
\(r^2(R+\varepsilon r)^2\) gives
\begin{equation}\label{eq:lagrange-cleared}
 GM_Sr^2+\varepsilon GM_E(R+\varepsilon r)^2
 -\omega^2r^2(R+\varepsilon r)^3=0.
\end{equation}

\begin{theorem}[The collinear equation is quintic]
\label{thm:lagrange-quintic}
If \(\omega\ne0\), the left side of~\eqref{eq:lagrange-cleared} is a polynomial
of degree five in \(r\), with leading coefficient \(-\varepsilon\omega^2\).
Equation~\eqref{eq:lagrange-force} is equivalent to this quintic under the two
denominator guards.
\end{theorem}

\begin{proof}
The last term has degree \(2+3=5\); the other terms have degree at most two.
Its leading monomial is
\(-\omega^2r^2(\varepsilon r)^3=-\varepsilon\omega^2r^5\).
Multiplication by the nonzero common denominator is reversible.  All statements
are polynomial identities and guarded field cancellation.
\end{proof}

Using Kepler's relation
\[
 \omega^2=\frac{G(M_S+M_E)}{R^3}.
\]
a convenient integral scaling of~\eqref{eq:lagrange-cleared} is
\begin{equation}\label{eq:lagrange-kepler}
 \begin{split}
 0={}&M_SR^3r^2+\varepsilon M_ER^3(R+\varepsilon r)^2\\
    &\quad-(M_S+M_E)r^2(R+\varepsilon r)^3.
 \end{split}
\end{equation}
Expanding with \(\varepsilon^2=1\) gives
\begin{equation}\label{eq:lagrange-expanded}
 \begin{split}
 0={}&-\varepsilon(M_S+M_E)r^5
      -3R(M_S+M_E)r^4
      -3\varepsilon R^2(M_S+M_E)r^3\\
    &+(\varepsilon-1)M_ER^3r^2
      +2M_ER^4r+\varepsilon M_ER^5.
 \end{split}
\end{equation}
An overall sign or the replacement \(r\mapsto-r\) changes the signs in a
coefficient table, which explains apparent convention differences among
references; the denominator-cleared equation~\eqref{eq:lagrange-cleared} is the
unambiguous certificate.

\subsection{The Hill-radius approximation}

Let
\[
 \mu=\frac{M_E}{M_S+M_E},
 \qquad s=\frac rR.
\]
For the upper-sign equation \(\varepsilon=1\), dividing
\eqref{eq:lagrange-kepler} by \((M_S+M_E)R^5\) yields the exact relation
\begin{equation}\label{eq:L2-dimensionless}
 \mu(1+2s)=s^3(3+3s+s^2).
\end{equation}
For small \(s\), the dominant balance is \(3s^3=\mu\), giving
\begin{equation}\label{eq:hill-approx}
 r\approx R\sqrt[3]{\frac{M_E}{3(M_S+M_E)}}
 \sim R\sqrt[3]{\frac{M_E}{3M_S}}.
\end{equation}
The lower-sign equation has the same leading balance.

This approximation can be certified without asymptotic notation.  Given a
rational upper bound \(s\le\delta<1/2\), the rational factor
\[
 A(s)=\frac{3+3s+s^2}{1+2s}.
\]
differs from \(3\) by an explicitly bounded rational multiple of \(s\).
Equation~\eqref{eq:L2-dimensionless} is \(\mu=A(s)s^3\).  Rational upper and
lower bounds on \(A\), followed by monotonicity of the cube, give an interval
around \((\mu/3)^{1/3}\).  Alternatively one evaluates the exact quintic at two
rational endpoints and invokes bisection.  Thus every requested relative error
in~\eqref{eq:hill-approx} has a primitive-recursive smallness threshold for
\(\mu\).

\subsection{A general certificate-oriented solver}

The preceding results suggest the following exact workflow for a rational
quintic.

\begin{enumerate}
\item Normalize and factor over \(\Q\).  Solve factors of degree at most four
      with the checked factorization identities.
\item For an irreducible quintic, gather modular factorization data and a
      discriminant.  This often identifies the Galois group.  If the group is
      solvable, construct Lagrange resolvents as in
      Theorem~\ref{thm:solvable-quintic-resolvents}.
\item If the group is nonsolvable, compute a principal and then a certified
      Bring--Jerrard transform.  Evaluate the required Bring branch by series,
      bisection, or analytic continuation, and invert the Tschirnhaus map by a
      subresultant certificate.
\item In every case, substitute the final candidates into the original
      polynomial and verify a complete factorization.  Numerical approximation
      chooses branches but never serves as the final proof.
\end{enumerate}

Every loop has a fixed degree bound and every final claim is an exact polynomial
or finite-table certificate.  This is the form of a quintic solver that is most
natural for \ipra{} and for proof-producing software.

\section{Lean formalization plan}\label{sec:lean}

The archive contains \code{Lean/BringQuinticBlueprint.lean} and a separate
roadmap.  The file is intentionally a blueprint rather than a claim of a new
completed build.  This section gives the precise path from the human proofs to
closed terms of the same form as the existing ProveIt derivations.

\subsection{Target theorem shape}

The existing derivation library represents a closed \ipra{} theorem as
\[
 ([]:\operatorname{List}(\mathrm{Fm}\;\mathrm{PTm}\;0))
 \;\vdash[\mathrm{PRAAx}]\;\varphi.
\]
A new theorem should have this type after all free variables are universally
closed.  The metatheoretic Lean proof may use the repository's quotient and
ring infrastructure, but the resulting object theorem contains only the
\ipra{} axioms.  The audit command should print the axioms of every exported
Lean theorem, exactly as in \code{IPRADerivations/Checks/Axioms.lean}.

The regular-real statements are theorem schemas.  A fixed primitive-recursive
coefficient or approximation function is introduced as a language symbol, and
its defining equations become part of \code{PRAAx}.  We never attempt to put a
Lean function variable \(\N\to\N\) inside first-order \ipra.

\subsection{Proposed module tree}

\begin{description}[style=nextline,leftmargin=0pt,labelindent=0pt,labelwidth=0pt,labelsep=0pt]
\item[\code{Algebra/Bring/Polynomial.lean}]
  Integer-difference terms for \(S_4\), the identity
  \eqref{eq:s4-sos}, strong monotonicity after importing signed order, the
  trinomial discriminant, scaling, and the differential-resolvent numerator.

\item[\code{Algebra/Bring/Rational.lean}]
  Reduced rational codes, exact order and midpoint, \(P_a(x)\), the bisection
  step, the interval invariant, \(\lceil\log_2\rceil\), and the rational
  approximation function \code{bringApprox}.

\item[\code{Algebra/Bring/RegularReal.lean}]
  The arithmetic predicate saying that the fixed approximation sequence is
  regular, compatibility of two parameter names, the Lipschitz and residual
  bounds, oddness, unboundedness, and the fifth-root asymptotic modulus.

\item[\code{Algebra/Bring/TruncatedSeries.lean}]
  Codes for vectors in \(\Q[X]/(X^{N+1})\), Cauchy product, composition,
  geometric inverse, finite residue, and the finite Lagrange theorem
  \eqref{eq:finite-lagrange}.

\item[\code{Algebra/Bring/Series.lean}]
  The Bring coefficient function, the finite identity
  \(Y_N^5+Y_N=X\bmod X^{N+1}\), the hypergeometric ratio, convergence
  majorants, and the branch-point resultant.

\item[\code{Algebra/Quintic/NormalForms.lean}]
  The depression identity, power-sum/Newton identities through degree five,
  the quadratic principal transform, generic Sylvester matrices, and the
  soundness of a Bring--Jerrard certificate.

\item[\code{Algebra/Quintic/Cyclotomic5.lean}]
  The quotient by \(1+W+W^2+W^3+W^4\), root-of-unity filters, the de Moivre
  product, the four-resolvent product, and the explicit examples.

\item[\code{Algebra/Quintic/FiniteGroups.lean}]
  Five-letter permutation codes, subgroup generation, commutators, the
  explicit proofs \([S_5,S_5]=A_5\), \([A_5,A_5]=A_5\), and the affine group
  derived series.

\item[\code{Algebra/Quintic/GaloisCertificates.lean}]
  Finite-field polynomial arithmetic, Frobenius traces, modular factorization
  certificates for the two concrete polynomials, and the finite decomposition
  group theorem needed to obtain cycle types.

\item[\code{Algebra/Quintic/RadicalTowers.lean}]
  Coded extension tables, roots of unity, Kummer kernels, solvable normal
  series, and Lagrange-resolvent reconstruction for a group contained in
  \(\operatorname{AGL}(1,5)\).
\end{description}

\subsection{First milestone: exact polynomial theorems}

This milestone uses only existing ring infrastructure and should be completed
first.  A representative object-theory definition is
\begin{lstlisting}[style=code,language={}]
-- schematic: variables are signed integer terms ZT n
abbrev S4 (x y : ZT n) : ZT n :=
  x^4 + x^3*y + x^2*y^2 + x*y^3 + y^4

def bringSOS : Fm PTm n :=
  ZT.eq (12 * S4 x y)
    (3 * (2*x^2 + x*y)^2 + (3*x*y + 2*y^2)^2 + 8*y^4)

theorem bring_sos :
  ([] : List (Fm PTm 0)) |-[PRAAx] univ n bringSOS := by
  refine closeUniv n (ZT.derive ?_)
  ipra_ring
\end{lstlisting}
The actual syntax must use the repository's \code{ZT.natLit}, multiplication
and power notation, as illustrated by
\code{Algebra/PolynomialFormulas.lean}.  The depression identity and both
cyclotomic products have the same shape: assumptions are introduced by
\code{impI}, converted into quotient-ring equalities by \code{ZT.toZ\_mem}, and
closed by \code{ipra\_ring} or \code{linear\_combination}.

The resultant identity can be proved in either of two ways:

\begin{enumerate}
\item instantiate the fixed \(9\times9\) determinant from
      \code{Core.LinearAlgebra} and normalize; or
\item prove the elimination implications in both directions using
      \(5X^4+p=0\) and \(4pX+5q=0\), together with a Bézout identity for the
      converse.
\end{enumerate}

The first route generalizes; the second produces a smaller term for the
trinomial.

\subsection{Second milestone: rational bisection}

Define total rational operations.  A denominator-zero input returns a fixed
default, and every theorem carries a positivity guard.  The bisection state is
a paired code \(\angles{L,U}\).  Its step is a primitive recursion using the
Boolean rational comparison
\[
 \chi(P_a((L+U)/2)\le0).
\]
The main induction theorem should state the three clauses of
Lemma~\ref{lem:bisection-invariant} as a single Boolean conjunction.  The
width clause is an equality after multiplying by \(2^n\), avoiding rational
exponents.

The regularity theorem for \code{bringApprox} follows from the nested-interval
invariant.  Soundness of the limit is represented arithmetically as:
for every precision \(k\), the residual of the \(k\)-th approximation is below
a displayed dyadic bound.  Equality to zero of the named real is the universal
closure of these residual bounds.  The Lipschitz theorem is first proved for
rational parameters and approximants, then transferred to the fixed regular
names by their moduli.

\subsection{Third milestone: truncated series}

A series vector can reuse the repository's pairing and bounded-sum modules.
Recommended definitions are
\begin{align*}
 \operatorname{coeff}(A,i)&=\text{the }i\text{-th decoded rational},\\
 \operatorname{mul}(A,B,n)_k
 &=\sum_{i\le k}A_iB_{k-i},\\
 \operatorname{compose}(A,B,n)
 &=\sum_{j\le n}A_jB^j\pmod{X^{n+1}}.
\end{align*}
The fixed-point proof uses the Boolean relation
\[
 \operatorname{Agree}(A,B,m)
 =\prod_{i<m}[A_i=B_i].
\]
One quantifier-free induction proves that one application of
\(A\mapsto X\Phi(A)\) increases the agreement bound by one.  A second induction
proves the finite residue substitution lemma.  These two results yield the
uniform finite Lagrange theorem.

For the Bring series, define
\begin{lstlisting}[style=code,language={}]
bringCoeff (n : Nat) : RatCode :=
  if h : n % 4 = 1 then
    let k := (n - 1) / 4
    (-1)^k * choose (5*k) k / (4*k+1)
  else 0
\end{lstlisting}
with the sign and division implemented on codes.  The object theorem is the
coefficient-vector equality, not a Lean equality of completed formal power
series.

\subsection{Fourth milestone: finite groups before Galois theory}

The finite group layer can be completed independently of field theory.  Encode
a permutation by a base-six numeral or a five-entry nested pair.  Verify
bijectivity by a five-term Boolean product.  Enumerate \(S_5\) by bounded search
through the \(5^5\) maps and filter permutations; this is primitive recursive,
although a hard-coded 120-entry table gives smaller proof terms.

The perfectness proof of \(A_5\) need not enumerate all commutators.  Formalize
the two explicit double transpositions from Lemma~\ref{lem:A5-perfect}, prove
their commutator is \((1\,2\,3)\), prove normality of the commutator subgroup,
and prove that 3-cycles generate \(A_5\).  The affine-group theorem can follow
the minimal-normal-subgroup proof of Theorem~\ref{thm:AGL-criterion}, with the
156-subgroup enumeration retained as a check theorem.

\subsection{Fifth milestone: the concrete \(S_5\) certificates}

Polynomial arithmetic modulo \(2,3,7\) is tiny.  The Frobenius remainder tables
and Bézout identities in Section~\ref{sec:abel-ruffini} can be inserted as
literal vectors and verified by \code{native\_decide} in the Lean metatheory;
a reflection lemma then turns an accepted vector computation into the
corresponding object-theory equality.

The nontrivial bridge is the factor-degree/Frobenius theorem.  A clean module
order is:

\begin{enumerate}
\item finite quotient rings and fields;
\item square-free factorization and Frobenius orbits;
\item coded finite free \(\Z\)-orders and reduction modulo a prime;
\item decomposition groups and the residue Frobenius lift;
\item cycle type on the roots.
\end{enumerate}

For the two concrete examples the splitting-field degree is at most 120, so all
objects can be bounded a priori.  The general theorem is nevertheless cleaner
than hard-coding one splitting field.

\subsection{Sixth milestone: radical towers and the generic theorem}

The finite radical-tower theorem uses basis tables and automorphism tables.
Kummer-step soundness is the equation
\((\sigma\alpha/\alpha)^n=1\).  A normal-series constructor appends the abelian
kernel at each step.  The generic \(S_5\) theorem requires a separate syntactic
ring of rational functions in five indeterminates and the constructive
symmetric-polynomial normalizer of Lemma~\ref{lem:symmetric-rational}.  It need
not quantify over arbitrary fields inside the object theory; it is a fixed
symbolic computation theorem about polynomial codes.

\subsection{Completion criteria}

A complete Lean delivery should satisfy all of the following.

\begin{enumerate}
\item \code{lake build IPRADerivations} succeeds with the new modules enabled.
\item No \code{sorry}, \code{admit}, or unproved axiom is introduced.
\item Every exported theorem has the closed \code{Prf PRAAx} type or is clearly
      marked as a metatheoretic helper.
\item \code{Checks/Axioms.lean} reports only the same Lean metatheory axioms
      already present in the repository, and no new object-theory axiom.
\item The Python certificate JSON is regenerated and its hashes match the
      archive manifest.
\item The formal statement audit points to the exact Lean theorem name for
      every item upgraded from \(\Outline\) to \(\Exact\) or
      \(\Effective\).
\end{enumerate}

The ordering above deliberately yields useful checked theorems early: the
polynomial identities, real bisection, series coefficients and finite groups do
not have to wait for the larger Galois infrastructure.

\section{Statement-by-statement audit}\label{sec:audit}

The following table records the mathematical coverage.  ``Exact'' means a
finite identity or finite certificate theorem whose \ipra{} proof is fully
given.  ``Effective'' includes explicit names and moduli.  ``Certificate''
means the verifier and soundness theorem are complete but a very large
universal witness or a page-specific substitution is not expanded here.
``Outline'' gives a concrete formalization route but not a completed proof of
the global source claim.

\small
\begin{longtable}{>{\raggedright\arraybackslash}p{0.27\textwidth}
                  >{\raggedright\arraybackslash}p{0.15\textwidth}
                  >{\raggedright\arraybackslash}p{0.50\textwidth}}
\caption{Coverage of mathematical statements from the three reference pages.}
\label{tab:audit}\\
\toprule
Statement & Status & Result and qualification\\
\midrule
\endfirsthead
\toprule
Statement & Status & Result and qualification\\
\midrule
\endhead

Unique real root of \(x^5+x+a\) & \Effective &
Theorem~\ref{thm:bring-exists}; exact rational bisection with a displayed
primitive-recursive modulus.\\

Oddness, monotone decrease and unboundedness of \(\BR\) & \Exact/\allowbreak\Effective &
Corollaries~\ref{cor:bring-odd-decreasing} and
\ref{cor:bring-unbounded}; based on the SOS identity
\eqref{eq:s4-sos}.\\

\(\BR(a)\sim-a^{1/5}\) & \Effective &
Theorem~\ref{thm:bring-asymptotic}; explicit absolute error
\(\le a^{-3/5}\) and relative error \(\le a^{-4/5}\) for \(a\ge1\).\\

Continuity/stability of the real branch & \Effective &
Theorem~\ref{thm:bring-lipschitz}; sharp bound
\(|\BR(a)-\BR(b)|\le|a-b|\) and residual certification.\\

Derivative of the real branch & \Effective &
Theorem~\ref{thm:bring-derivative};
\(\BR'(a)=-(1+5\BR(a)^4)^{-1}\) with local moduli.\\

Complex Bring radical is five-valued & \Exact &
The polynomial has degree five; branches are represented in a formal splitting
algebra.  Existence as five analytic values away from the discriminant is the
effective local inverse construction.\\

Four branch points and need for cuts & \Exact/\allowbreak\Effective &
Equation~\eqref{eq:branch-equation}, Theorem~\ref{thm:trinomial-discriminant},
and the simple-critical-point Puiseux argument.  Global monodromy is represented
by finite path certificates.\\

Taylor series of \(\BR\) & \Exact &
Theorem~\ref{thm:bring-coeff}; finite Lagrange inversion in every truncation.\\

Hypergeometric form of \(\BR\) & \Exact/\allowbreak\Effective &
Theorem~\ref{thm:bring-hypergeom}; exact coefficient recurrence and effective
tails.\\

Radius \(4/(5\sqrt[4]5)\) & \Effective &
Theorem~\ref{thm:bring-radius}; constructive ratio modulus, boundary
convergence, and matching critical values.\\

General trinomial discriminant & \Exact &
\(\Disc(X^5+pX+q)=256p^5+3125q^4\),
Theorem~\ref{thm:trinomial-discriminant}.\\

Cockle--Harley fourth-order ODE & \Exact/\allowbreak\Effective &
Theorem~\ref{thm:differential-resolvent}; rational derivative elimination,
cleared by \((5t^4-1)^7\).\\

All five local branches of \(t^5-t+\rho\) & \Effective &
Finite fixed-point construction around \(0,\pm1,\pm i\), followed by the common
differential resolvent.  The exact linear combinations in MathWorld's printed
hypergeometric basis are not separately normalized.\\

General quintic depression & \Exact &
Theorem~\ref{thm:depress-quintic}, including all four coefficient formulas.\\

Derivative of a quintic is quartic & \Exact &
Formal differentiation; real-root isolation is a bounded Sturm or
critical-point algorithm.\\

Quadratic reduction to principal form & \Exact &
Theorem~\ref{thm:principal-reduction}; explicit \(\alpha,\beta\) equations from
power sums.\\

Universal quartic Tschirnhaus reduction to Bring--Jerrard form & \Conditional &
Theorem~\ref{thm:bring-certificate-sound} gives a compact resultant verifier and
complete soundness proof.  The classical universal witness, obtained through a
cubic and quadratic, is not expanded into its multimegabyte symbolic form.\\

Scaling \(x^5+px+q\) to a Bring equation & \Exact &
Theorem~\ref{thm:bring-scaling}; branch and fourth-root compatibility are made
explicit.\\

Inverting the Tschirnhaus map by equations of degree at most four & \Conditional &
Certified subresultant/gcd interface following~\eqref{eq:inverse-system};
degenerate degrees are recorded rather than ignored.\\

Brioschi normal form & \Conditional &
Any proposed rational transformation is checked after denominator clearing.
The global polyhedral formulas for its parameters remain with the icosahedral
outline.\\

General quintic not solvable by radicals & \Exact &
Theorem~\ref{thm:abel-ruffini}; generic group \(S_5\), radical towers solvable,
and explicit finite commutator proofs.\\

Solvability criterion for irreducible quintics & \Exact &
Theorem~\ref{thm:AGL-criterion}; solvable transitive subgroups are conjugate
into \(\operatorname{AGL}(1,5)\).  The independent enumeration checks all 156
subgroups.\\

\(x^5-x+1\) is not solvable by radicals & \Exact/\allowbreak\Conditional &
Theorem~\ref{thm:x5-x+1-S5}; exact modular factors and Frobenius traces.  The
factor-degree/Frobenius bridge is supplied as a finite Galois-certificate
module plan rather than a completed new Lean file.\\

A concrete Bring value is not radical & \Exact/\allowbreak\Conditional &
Theorem~\ref{thm:BR3-not-radical}: \(x^5+x+3\) has group \(S_5\), so its real
root \(\BR(3)\) is not in a radical extension of \(\Q\).  Same Galois-module
qualification as the preceding row.\\

Reducible example \(x^5-x^4-x+1\) & \Exact &
Factorization~\eqref{eq:reducible-example}.\\

de Moivre quintic and auxiliary quadratic & \Exact &
Theorem~\ref{thm:de-moivre-product} and
Corollary~\ref{cor:de-moivre-roots}.\\

General four-resolvent expansion & \Exact &
Theorem~\ref{thm:four-resolvent}; direct cyclotomic quotient identity matching
MathWorld's coefficient invariants.\\

Spearman--Williams radical root expressions & \Conditional &
Corollary~\ref{cor:four-resolvent-specialization} is the complete root verifier.
The long substitution showing that the page's specific \(u_i\) satisfy the
four invariant equations is left as a generated certificate.\\

Spearman--Williams ``if and only if'' rational parameterization & \Outline &
Reduced to rational invariant theory plus the affine-group criterion; no claim
of a completed normalization of all published parameter formulas.\\

Runge/Cadenhad/Young parameterization & \Conditional &
The displayed coefficients are exact rational maps and can feed the universal
resolvent verifier.  Completeness of the parameterization is not reproved.\\

Explicit root involving \(\sqrt[5]2\) & \Exact &
Equation~\eqref{eq:explicit-example-one}; exact remainder modulo \(t^5-2\).\\

Paxton Young example with \(\sqrt[5]2\) & \Exact &
Equation~\eqref{eq:young-example}.\\

Root \(2\cos(2\pi/11)\) of a quintic & \Exact &
Equation~\eqref{eq:cos11-example}; reduction modulo the eleventh cyclotomic
polynomial.\\

The source's quartic-resolvent formula for the root of \(x^5-5x+12\) & \Conditional &
It fits Theorem~\ref{thm:solvable-quintic-resolvents}; the particular four
fifth-root compatibility relations are not expanded in this packet.\\

The five listed solvable \(x^5+ax^2+b\) shapes & \Outline &
Each instance is decidable by a Galois/resolvent certificate, but the historical
classification ``exactly five up to scaling'' is not reproved here.\\

Kondo--Brumer and other parameterized solvable families & \Conditional &
A supplied group/resolvent certificate is checkable.  Universal family
solvability is a future generated-certificate target.\\

Casus irreducibilis for quintics & \Outline &
Ordered radical-tower argument in Section~\ref{sec:solvable}; individual towers
are checkable, but the full ordered-field Lean module is not complete.\\

Hermite theta-function solution & \Outline &
Theorem~\ref{thm:hermite-interface} gives exact data and tail requirements.
Global modular equations and branch choices are not completed.\\

Klein icosahedral/hypergeometric solution & \Outline &
Local hypergeometric and polynomial verification are finite; global monodromy
and invariant normalization remain.\\

Hypergeometric solution for \(x^p+bx^q+c\) & \Effective{} schema &
Finite Lagrange inversion and coefficient recurrences for every fixed
\((p,q)\); global branch formulas require case-specific continuation.\\

Doyle--McMullen iteration & \Conditional &
Theorem~\ref{thm:certified-iteration}; a trapping disk and contraction bound
replace the nonformal instruction to use a random start.\\

Quintic from collinear Lagrange points & \Exact &
Theorem~\ref{thm:lagrange-quintic} and expansion
\eqref{eq:lagrange-expanded}.\\

Hill-radius approximation & \Effective &
Exact dimensionless equation~\eqref{eq:L2-dimensionless}; rational interval
bounds give any requested asymptotic accuracy.\\

Historical attributions and dates & Not arithmetic &
Recorded in the source bibliography but not presented as theorems of \ipra.\\

\bottomrule
\end{longtable}
\normalsize

\subsection{Net result}

The largest fully developed fragment is the one most naturally suited to
\ipra: polynomial identities, finite series coefficients, effective real
roots, finite groups, radical-tower certificates, and exact numerical
verification.  The remaining gaps are concentrated in global special-function
identities and the size of universal symbolic transformations, not in the
constructive treatment of the Bring radical itself.

\clearpage
\begin{thebibliography}{99}

\bibitem{wikiBring}
Wikipedia contributors,
\emph{Bring radical},
\url{https://en.wikipedia.org/wiki/Bring_radical},
accessed 16 November 2026.

\bibitem{wikiQuintic}
Wikipedia contributors,
\emph{Quintic equation},
\url{https://en.wikipedia.org/wiki/Quintic_equation},
accessed 16 November 2026.

\bibitem{mathworldQuintic}
Eric W. Weisstein,
\emph{Quintic Equation}, MathWorld---A Wolfram Resource,
\url{https://mathworld.wolfram.com/QuinticEquation.html},
accessed 16 November 2026.

\bibitem{proveit}
Vladimir Reshetnikov,
\emph{ProveIt}, GitHub repository,
\url{https://github.com/VladimirReshetnikov/ProveIt},
repository state headed by commit
\code{b73ff58ad10e9c446fa997777c5f502844249929}.

\bibitem{proveitPRA}
ProveIt,
\emph{Logic/PrimitiveRecursiveArithmetic},
\url{https://github.com/VladimirReshetnikov/ProveIt/tree/main/Logic/PrimitiveRecursiveArithmetic}.

\bibitem{proveitDerivations}
ProveIt,
\emph{IPRADerivations}, including the checked
\code{Algebra/PolynomialFormulas.lean},
\url{https://github.com/VladimirReshetnikov/ProveIt/tree/main/Logic/PrimitiveRecursiveArithmetic/Derivations}.

\bibitem{proveitLagrange}
ProveIt,
\emph{Lagrange inversion in intuitionistic primitive recursive arithmetic},
\url{https://github.com/VladimirReshetnikov/ProveIt/tree/main/Logic/PrimitiveRecursiveArithmetic/docs/proofs/algebra/lagrange_inversion_ipra}.

\bibitem{proveitAbel}
ProveIt,
\emph{Polynomial formulas and Abel--Ruffini in iPRA},
\url{https://github.com/VladimirReshetnikov/ProveIt/tree/main/Logic/PrimitiveRecursiveArithmetic/docs/proofs/algebra/ipra_polynomial_formulas_abel_ruffini}.

\bibitem{Abel}
Niels Henrik Abel,
\emph{Mémoire sur les équations algébriques, où on démontre l'impossibilité de la résolution de l'équation générale du cinquième degré},
1824/1826.

\bibitem{Galois}
Évariste Galois,
\emph{Mémoire sur les conditions de résolubilité des équations par radicaux},
1831/1846.

\bibitem{Hermite}
Charles Hermite,
\emph{Sur la résolution de l'équation du cinquième degré},
Comptes rendus de l'Académie des Sciences 46 (1858), 508--515.

\bibitem{Klein}
Felix Klein,
\emph{Lectures on the Icosahedron and the Solution of Equations of the Fifth Degree},
1888; English translation, Trübner \& Co.

\bibitem{SpearmanWilliams}
Blair K. Spearman and Kenneth S. Williams,
\emph{Characterization of solvable quintics \(x^5+ax+b\)},
American Mathematical Monthly 101 (1994), 986--992.

\bibitem{Runge}
Carl Runge,
\emph{Über die auflösbaren Gleichungen von der Form \(x^5+ux+v=0\)},
Acta Mathematica 7 (1885), 173--186.

\bibitem{AdamchikJeffrey}
Victor S. Adamchik and David J. Jeffrey,
\emph{Polynomial transformations of Tschirnhaus, Bring and Jerrard},
ACM SIGSAM Bulletin 37(3) (2003), 90--94.

\bibitem{Dummit}
David S. Dummit,
\emph{Solving solvable quintics}, expository manuscript,
archived at
\url{https://web.archive.org/web/20120307005522/http://www.emba.uvm.edu/~dummit/quintics.pdf}.

\bibitem{DoyleMcMullen}
Peter Doyle and Curt McMullen,
\emph{Solving the quintic by iteration},
Acta Mathematica 163 (1989), 151--180.

\bibitem{King}
R. Bruce King,
\emph{Beyond the Quartic Equation},
Birkhäuser, 1996.

\end{thebibliography}


\end{document}
